\chapter{Independent prime divisors}\label{ch:oc-divisors}\label{ch:additive}\label{ch:sd}\label{ch:erdoskac}

The prime factors of a typical integer are nearly independent. This chapter proves the Tur\'an--Kubilius inequality, the Kubilius model, and the Selberg--Delange theorem that Volume~11 quotes from this volume. The Erd\H{o}s--Kac theorem is a corollary of the model, set beside the proof in Volume~6, Chapter~9. It is not founded a second time.

\section{Additive functions and the Kubilius model}

\begin{center}\itshape
Pick an integer at random. Is it abundant? How close is $\varphi(n)$ to $n$?\\
The primes dividing $n$ behave like independent coins,\\
and an additive function is the sum of what they show.
\end{center}

In 1928 Isaac Schoenberg, a young Romanian mathematician then working in Germany, proved a curious fact about Euler's function: the proportion of integers $n\le x$ with $\frac{\varphi(n)}n\le z$ tends, for every $z$, to a limit $F(z)$, and $F$ is continuous. Five years later Harold Davenport proved the same kind of statement for $\frac{\sigma(n)}n$, the sum of the divisors of $n$ divided by $n$. As a consequence the \emph{abundant numbers}, those with $\sigma(n)\ge2n$, have a natural density, a question that went back to the Greeks' fascination with perfect numbers. In the same year Felix Behrend showed that the density lies between $0.241$ and $0.314$. It is now known to be $0.24761\ldots$, which means that roughly one integer in four is abundant.

These were hard-won theorems, each proved with its own machinery. In 1939 Paul Erd\H{o}s and Aurel Wintner, in a nine-page paper in the \emph{American Journal of Mathematics}, explained them all at once. Write $\log\frac{\varphi(n)}n=\sum_{p\mid n}\log(1-\frac1p)$. This is an \emph{additive} function: a sum over the primes dividing $n$. Erd\H{o}s and Wintner gave a criterion for an additive function to have a limiting distribution, and the criterion is the three-series theorem that Kolmogorov had proved in 1928 for sums of independent random variables. The title of their paper says why: \emph{Additive arithmetical functions and statistical independence}.

The independence is a theorem and not just a heuristic. Twenty years later Jonas Kubilius in Vilnius made it precise: the exponents of the small primes in a random integer $n\le x$ are, up to a tiny error in total variation, independent geometric random variables. Kubilius became rector of Vilnius University in 1958 and led it for more than thirty years, until 1990, while building the Lithuanian school of probability and number theory. His monograph, translated by the American Mathematical Society in 1964 as \emph{Probabilistic Methods in the Theory of Numbers}, gave the subject its name.

This chapter builds the model and draws the consequences: the Tur\'an--Kubilius inequality for every additive function, the Kubilius model in total variation, the Erd\H{o}s--Wintner theorem, and two theorems of probability, of L\'evy and of Jessen and Wintner, which decide whether the limit law is continuous. Euler's function and the abundant numbers come out as corollaries.

\begin{reminder}[from Volumes~4, 6, 11 and~14]
\emph{Tur\'an--Kubilius} (Volume~6, Theorem~9.1). There is an absolute $C_0$ such that for every real function $f$ on the primes and $x\ge3$, the strongly additive $f(n)=\sum_{p\mid n}f(p)$ satisfies $\sum_{n\le x}(f(n)-A_0(x))^2\le C_0xB_0(x)^2$, where $A_0(x)=\sum_{p\le x}\frac{f(p)}p$ and $B_0(x)^2=\sum_{p\le x}\frac{f(p)^2}p$.
\emph{Rough numbers} (Volume~14, Corollary~3.5). For $x\ge4$ and $2\le z\le\sqrt x$, the number $\Phi(x,z)$ of $n\le x$ with no prime factor $\le z$ is $x\prod_{p\le z}(1-\frac1p)(1+O(e^{-u/2}))+O(\sqrt x\log x)$, with $u=\frac{\log x}{\log z}$.
\emph{Three-series theorem} (Volume~11, Theorem~2.5). For independent $X_p$ and $c>0$, $\sum X_p$ converges almost surely if and only if $\sum P(|X_p|>c)$, $\sum\mathbb E[X_p\mathbf 1_{|X_p|\le c}]$ and $\sum\operatorname{Var}[X_p\mathbf 1_{|X_p|\le c}]$ converge.
\emph{Kolmogorov's zero--one law} (Volume~11, Theorem~1.4) and the \emph{Borel--Cantelli lemmas} (Volume~4, Chapter~4).
\emph{Mertens} (Volume~6, Theorems~1.3--1.6). $\sum_{p\le y}\frac{\log p}p=\log y+O(1)$, $\sum_{p\le y}\frac1p=\log\log y+O(1)$ and $\prod_{p\le y}(1-\frac1p)\asymp\frac1{\log y}$.
\end{reminder}

\section{Additive functions}

\begin{definition}\label{def:additive}
A function $f:\N_{\ge1}\to\C$ is \emph{additive} if $f(mn)=f(m)+f(n)$ whenever $\gcd(m,n)=1$, so that $f(n)=\sum_{p^k\parallel n}f(p^k)$, where $p^k\parallel n$ means that $p^k$ divides $n$ and $p^{k+1}$ does not. It is \emph{strongly additive} if moreover $f(p^k)=f(p)$ for all $k\ge1$, and \emph{completely additive} if $f(mn)=f(m)+f(n)$ for all $m,n$.
\end{definition}

The number $\omega(n)$ of distinct prime factors is strongly additive, $\Omega(n)$ (with multiplicity) and $\log n$ are completely additive, and $\log\frac{\varphi(n)}n=\sum_{p\mid n}\log(1-\frac1p)$ is strongly additive. If $g$ is multiplicative and positive, $\log g$ is additive; so $\log\frac{\sigma(n)}n$ is additive, with $\log\frac{\sigma(p^k)}{p^k}=\log\frac{1-p^{-k-1}}{1-p^{-1}}$.

The object of study is the \emph{distribution} of $f$: for $n$ chosen uniformly in $\{1,\dots,\lfloor x\rfloor\}$, the law of the random variable $f(n)$. We write $P_x$ and $\mathbb E_x$ for probabilities and expectations under this uniform choice, so that $P_x(f(n)\le z)=\frac1{\lfloor x\rfloor}\#\{n\le x:f(n)\le z\}$. We say that $f$ has a \emph{limiting distribution} if $f(n)$ converges in law as $x\to\infty$, that is, if $P_x(f(n)\le z)$ tends to $F(z)$ at every continuity point $z$ of some distribution function $F$.

\begin{think}
Which of $\omega(n)$, $\log\frac{\varphi(n)}n$ and $\log\frac{\sigma(n)}n$ would you expect to have a limiting distribution without any renormalisation? Think of the size of the contribution of each prime, and of how often that prime divides $n$.
\end{think}

The heuristic is simple. A prime $p$ divides a random integer with probability $\frac1p$, and does so ``independently'' of the other primes. For $\omega$ each prime contributes $1$ with probability $\frac1p$, and $\sum\frac1p$ diverges: $\omega(n)$ grows like $\log\log n$ and must be centred and scaled, which is the Erd\H{o}s--Kac theorem of Volume~6 and of this chapter. For $\log\frac{\varphi(n)}n$ the prime $p$ contributes about $-\frac1p$ with probability $\frac1p$, and $\sum\frac1{p^2}$ converges: the sum should converge without any normalisation. The rest of the chapter makes this precise.

\section{The Tur\'an--Kubilius inequality}

Volume~6 proved the inequality for strongly additive functions. Every additive function differs from a strongly additive one only at the prime powers $p^k$ with $k\ge2$, which are rare, and that is the whole proof of the general case.

\begin{theorem}[Tur\'an--Kubilius]\label{thm:tkgeneral}
There is an absolute constant $C_1$ such that for every additive $f:\N\to\C$ and every $x\ge3$,
\[
\sum_{n\le x}|f(n)-A(x)|^2\le C_1xB(x)^2,\qquad A(x)=\sum_{p^k\le x}\frac{f(p^k)}{p^k}\Bigl(1-\frac1p\Bigr),\qq B(x)^2=\sum_{p^k\le x}\frac{|f(p^k)|^2}{p^k}.
\]
One can take $C_1=3C_0+36$.
\end{theorem}

The centring $A(x)$ is the expectation of the \emph{model}: if $G_p$ is a geometric variable with $P(G_p=k)=(1-\frac1p)p^{-k}$, then $\mathbb E f(p^{G_p})=\sum_{k\ge1}f(p^k)(1-\frac1p)p^{-k}$, with $f(1)=0$.

\begin{proof}
Since $|w|^2=(\operatorname{Re}w)^2+(\operatorname{Im}w)^2$ and the real and imaginary parts of $f$ are additive with $A$ and $B^2$ splitting accordingly ($B^2=B_{\operatorname{Re}}^2+B_{\operatorname{Im}}^2$), it suffices to treat real $f$ with the constant $3C_0+36$.

For $n\le x$ write $f(n)=g(n)+h(n)$ with
\[
g(n)=\sum_{p\mid n}f(p),\qquad h(n)=\sum_{\substack{p^k\parallel n\\k\ge2}}c(p^k),\qquad c(p^k)=f(p^k)-f(p).
\]
The function $g$ is strongly additive, and Volume~6, Theorem~9.1, gives $\sum_{n\le x}(g(n)-A_0)^2\le C_0xB_0^2\le C_0xB^2$, with $A_0=\sum_{p\le x}\frac{f(p)}p$.

\emph{The function $h$.} Expanding the square, the diagonal terms are $\sum_{p^k}c(p^k)^2\#\{n\le x:p^k\parallel n\}\le x\sum_{p^k}\frac{c(p^k)^2}{p^k}$. In the off-diagonal terms $p^k\parallel n$ and $q^l\parallel n$ with $p^k\ne q^l$ force $p\ne q$, so $p^kq^l\mid n$ and there are at most $\frac x{p^kq^l}$ such $n$. Hence, with all sums over $p^k\le x$ with $k\ge2$, and by Cauchy--Schwarz,
\[
\sum_{n\le x}h(n)^2\le x\sum\frac{c(p^k)^2}{p^k}+x\Bigl(\sum\frac{|c(p^k)|}{p^k}\Bigr)^2\le x\,H\Bigl(1+\sum_p\sum_{k\ge2}\frac1{p^k}\Bigr)\le2xH,
\]
where $H=\sum\frac{c(p^k)^2}{p^k}$, because $\sum_p\sum_{k\ge2}p^{-k}=\sum_p\frac1{p(p-1)}<\sum_{m\ge2}\frac1{m(m-1)}=1$. Now $c(p^k)^2\le2f(p^k)^2+2f(p)^2$, and $\sum_{k\ge2}p^{-k}=\frac1{p(p-1)}\le\frac1p$, so
\[
H\le2\sum_{\substack{p^k\le x\\k\ge2}}\frac{f(p^k)^2}{p^k}+2\sum_{p\le x}\frac{f(p)^2}p\le4B^2,\qquad\text{and}\qquad\sum_{n\le x}h(n)^2\le8xB^2.
\]

\emph{The centring.} Since $\frac{f(p)}p-\frac{f(p)}p(1-\frac1p)=\frac{f(p)}{p^2}$,
\[
A_0-A=\sum_{p\le x}\frac{f(p)}{p^2}-\sum_{\substack{p^k\le x\\k\ge2}}\frac{f(p^k)}{p^k}\Bigl(1-\frac1p\Bigr),
\]
and by Cauchy--Schwarz each sum is at most $B$ in absolute value: $|\sum_p\frac{f(p)}{p^2}|\le(\sum_p\frac{f(p)^2}p)^{1/2}(\sum_p\frac1{p^3})^{1/2}\le B$ and $|\sum\frac{f(p^k)}{p^k}|\le B(\sum_p\sum_{k\ge2}p^{-k})^{1/2}\le B$. So $|A_0-A|\le2B$.

Finally $f(n)-A=(g(n)-A_0)+h(n)+(A_0-A)$ and $(a+b+c)^2\le3(a^2+b^2+c^2)$ give $\sum_{n\le x}(f(n)-A)^2\le3(C_0+8+4)xB^2$.
\end{proof}

In the language of probability: $\operatorname{Var}_xf(n)\le\mathbb E_x|f(n)-A(x)|^2\le C_1B(x)^2(1+O(\frac1x))$, and $B(x)^2$ is, up to bounded factors, the variance of the model $\sum_{p\le x}f(p^{G_p})$. The dependence between the primes costs at most a constant factor. Chebyshev's inequality turns this into a law of large numbers: if $B(x)\to\infty$, then $\frac{f(n)-A(x)}{B(x)\,\psi(x)}\to0$ in probability for every $\psi(x)\to\infty$. For $\omega$ this is the theorem of Hardy and Ramanujan (Volume~6, Theorem~9.2).

\section{The Kubilius model}

The inequality compares only two moments. Kubilius's insight was that much more is true: for the primes up to a small power of $x$, the whole joint law of the exponents is close to the independent model, in the strongest sense available.

For $y\ge2$ and $n\ge1$ let $v_y(n)=(\nu_p(n))_{p\le y}$ be the vector of exponents of the primes $p\le y$ in $n$. Let $(G_p)_p$ be independent with $P(G_p=k)=(1-\frac1p)p^{-k}$ for $k\ge0$, and $G^{(y)}=(G_p)_{p\le y}$. Recall that the \emph{total variation distance} between two laws $\mu,\nu$ on a countable set is $d_{TV}(\mu,\nu)=\sup_E|\mu(E)-\nu(E)|=\frac12\sum_a|\mu(a)-\nu(a)|$.

\begin{theorem}[Kubilius]\label{thm:kubilius}
There is an absolute constant $C_2$ such that for $x\ge16$ and $2\le y\le x^{1/4}$, with $u=\frac{\log x}{\log y}$,
\[
d_{TV}\bigl(\operatorname{law}_x v_y(n),\ \operatorname{law}G^{(y)}\bigr)\le C_2\bigl(e^{-u/4}+x^{-1/4}\log x\bigr).
\]
\end{theorem}

\begin{proof}
A vector $\mathbf k=(k_p)_{p\le y}$ with finitely many non-zero entries corresponds to a $y$-\emph{smooth} integer $m=\prod_pp^{k_p}$, one all of whose prime factors are $\le y$. Write $V=\prod_{p\le y}(1-\frac1p)$, so $P(G^{(y)}=\mathbf k)=\frac Vm$. On the other side, $v_y(n)=\mathbf k$ means $n=mn'$ with $n'$ free of prime factors $\le y$, so
\[
P_x(v_y(n)=\mathbf k)=\frac{\Phi(x/m,y)}{\lfloor x\rfloor}.
\]
\emph{Small $m$.} Let $m\le\sqrt x$. Then $X=\frac xm\ge\sqrt x\ge y^2$, and Volume~14, Corollary~3.5, applies with $u_X=\frac{\log X}{\log y}\ge\frac u2$:
\[
\Bigl|\frac{\Phi(x/m,y)}x-\frac Vm\Bigr|\ll\frac Vme^{-u/4}+\frac{\sqrt{x/m}\log x}x.
\]
Summing over $y$-smooth $m\le\sqrt x$, and using $\sum_{m\ y\text{-smooth}}\frac1m=\prod_{p\le y}(1-\frac1p)^{-1}=\frac1V$ and $\sum_{m\le\sqrt x}m^{-1/2}\le2x^{1/4}$,
\[
\Sigma_1:=\sum_{m\le\sqrt x}\Bigl|\frac{\Phi(x/m,y)}x-\frac Vm\Bigr|\ll e^{-u/4}+x^{-1/4}\log x.
\]
\emph{Large $m$, model side (Rankin's trick).} Let $\delta=\frac1{\log y}$. For $m>\sqrt x$, $1\le(\frac m{\sqrt x})^\delta$, so
\[
\Sigma_2:=\sum_{\substack{m>\sqrt x\\y\text{-smooth}}}\frac Vm\le Vx^{-\delta/2}\sum_{y\text{-smooth}}\frac1{m^{1-\delta}}=Ve^{-u/2}\prod_{p\le y}\Bigl(1-\frac1{p^{1-\delta}}\Bigr)^{-1}.
\]
For $p\le y$, $p^\delta=e^{\log p/\log y}\le1+(e-1)\frac{\log p}{\log y}$ by convexity, so $\sum_{p\le y}p^{\delta-1}\le\sum_{p\le y}\frac1p+\frac{e-1}{\log y}\sum_{p\le y}\frac{\log p}p=\log\log y+O(1)$, and $\sum_pp^{2\delta-2}\le\sum_pe^2p^{-2}=O(1)$. With $\log(1-t)^{-1}\le t+t^2$ for $0\le t\le\frac12$ (and the factor at $p=2$ bounded), the product is $\ll\log y$, and $V\ll\frac1{\log y}$. Hence $\Sigma_2\ll e^{-u/2}$.

\emph{Large $m$, integer side.} The probabilities $P_x(v_y(n)=\mathbf k)$ sum to $1$, and so do the $\frac Vm$. Since $\lfloor x\rfloor\le x$,
\[
\sum_{m>\sqrt x}\frac{\Phi(x/m,y)}{\lfloor x\rfloor}=1-\sum_{m\le\sqrt x}\frac{\Phi(x/m,y)}{\lfloor x\rfloor}\le1-\sum_{m\le\sqrt x}\frac{\Phi(x/m,y)}x\le1-\sum_{m\le\sqrt x}\frac Vm+\Sigma_1=\Sigma_2+\Sigma_1.
\]
\emph{Conclusion.} Since $\sum_{m\le\sqrt x}\Phi(x/m,y)\le x$, replacing $\frac1x$ by $\frac1{\lfloor x\rfloor}$ in $\Sigma_1$ costs at most $\frac x{\lfloor x\rfloor}-1\le\frac2x$. So
\[
2d_{TV}=\sum_{\mathbf k}\bigl|P_x(v_y(n)=\mathbf k)-P(G^{(y)}=\mathbf k)\bigr|\le\Bigl(\Sigma_1+\frac2x\Bigr)+(\Sigma_1+\Sigma_2)+\Sigma_2\ll e^{-u/4}+x^{-1/4}\log x.
\]
\end{proof}

\begin{think}
With $y=x^{1/u}$, the theorem is useful only when $u$ is large. What goes wrong for $y=\sqrt x$? (Count the $n\le x$ that are prime, and compare with the model's probability that $G_p=0$ for every $p\le\sqrt x$.)
\end{think}

For $y=\sqrt x$ the integer side has $P_x(v_y(n)=\mathbf0)=\frac{1+\pi(x)-\pi(\sqrt x)}{\lfloor x\rfloor}\sim\frac1{\log x}$, while the model gives $V\sim\frac{2e^{-\gamma}}{\log x}$: the two differ by the factor $\frac{e^\gamma}2\approx0.89$ already met in Volume~14, Figure~3.1. A large prime factor uses up most of the room below $x$, and that constraint is invisible to independent coins. The model is exact for the primes up to $x^{o(1)}$, and only for those.

Since $f(n)$ restricted to the small primes is a function of $v_y(n)$, the theorem transfers at once to every additive function. Let
\[
f_y(n)=\sum_{\substack{p^k\parallel n\\p\le y}}f(p^k),\qquad Z_y=\sum_{p\le y}f(p^{G_p})\quad(f(1)=0).
\]

\begin{corollary}\label{cor:kubiliusf}
For every additive $f$, $x\ge16$ and $2\le y\le x^{1/4}$, and every Borel set $E\subseteq\C$, $|P_x(f_y(n)\in E)-P(Z_y\in E)|\le C_2(e^{-u/4}+x^{-1/4}\log x)$.
\end{corollary}
\begin{proof}
$f_y(n)=\Lambda(v_y(n))$ and $Z_y=\Lambda(G^{(y)})$ for the same map $\Lambda(\mathbf k)=\sum_pf(p^{k_p})$, and total variation distance can only decrease under a map: $\{f_y(n)\in E\}=\{v_y(n)\in\Lambda^{-1}(E)\}$.
\end{proof}

\section{The Erd\H{o}s--Wintner theorem}

For a real additive $f$ and a prime $p$ write $f^*(p)=f(p)$ if $|f(p)|\le1$ and $f^*(p)=0$ otherwise.

\begin{theorem}[Erd\H{o}s and Wintner, 1939]\label{thm:erdoswintner}
Let $f$ be a real additive function such that the three series
\[
\sum_{|f(p)|>1}\frac1p,\qquad\sum_p\frac{f^*(p)}p,\qquad\sum_p\frac{f^*(p)^2}p
\]
converge. Then $f$ has a limiting distribution: $f(n)$, for $n$ uniform in $[1,x]$, converges in law as $x\to\infty$ to
\[
Z=\sum_pf(p^{G_p}),
\]
a series of independent variables that converges almost surely. Its characteristic function is
\[
\mathbb Ee^{itZ}=\prod_p\Bigl(1-\frac1p\Bigr)\Bigl(1+\sum_{k\ge1}\frac{e^{itf(p^k)}}{p^k}\Bigr).
\]
\end{theorem}

\begin{proof}
\emph{The model converges.} Let $X_p=f(p^{G_p})$ and apply the three-series theorem with $c=1$. Since $P(G_p\ge2)=p^{-2}$,
\[
\begin{gathered}P(|X_p|>1)\le\frac1{p^2}+\frac{\mathbf 1_{|f(p)|>1}}p,\\\Bigl|\mathbb E[X_p\mathbf 1_{|X_p|\le1}]-f^*(p)\frac{p-1}{p^2}\Bigr|\le\frac1{p^2},\qquad\mathbb E[X_p^2\mathbf 1_{|X_p|\le1}]\le\frac{f^*(p)^2}p+\frac1{p^2}.\end{gathered}
\]
The three hypotheses, with $\sum_p\frac1{p^2}<\infty$ and $|\frac{f^*(p)}{p^2}|\le\frac1{p^2}$, make the three series of Volume~11, Theorem~2.5, converge (the variance is at most the second moment). So $Z_y\to Z$ almost surely as $y\to\infty$, hence in law. Since $Z$ is the almost sure limit of independent sums, $\mathbb Ee^{itZ}=\lim_y\prod_{p\le y}\mathbb Ee^{itX_p}$ by dominated convergence, which is the product in the statement.

\emph{The small primes.} Let $y=x^{1/\log\log x}$, so $u=\log\log x\to\infty$ and $y\to\infty$. By Corollary~\ref{cor:kubiliusf} the laws of $f_y(n)$ and $Z_y$ are $o(1)$ apart in total variation. Since $Z_y\to Z$ in law, so does $f_y(n)$: for every bounded continuous $\phi$, $|\mathbb E_x\phi(f_y(n))-\mathbb E\phi(Z_y)|\le2\|\phi\|_\infty d_{TV}\to0$.

\emph{The large primes.} Let $R(n)=f(n)-f_y(n)=\sum_{p^k\parallel n,\,p>y}f(p^k)$. We show that $R(n)\to0$ in probability; Slutsky's lemma (Volume~4, Exercise~6.1) then gives $f(n)=f_y(n)+R(n)\to Z$ in law. Let $\tilde f$ be the strongly additive function with $\tilde f(p)=f^*(p)$ for $y<p\le x$ and $\tilde f(p)=0$ otherwise. Then $R(n)=\tilde f(n)$ unless $p^2\mid n$ for some $p>y$, or $p\mid n$ for some $p>y$ with $|f(p)|>1$. So
\[
P_x(R(n)\ne\tilde f(n))\le\sum_{p>y}\frac1{p^2}+\sum_{\substack{p>y\\|f(p)|>1}}\frac1p+O\Bigl(\frac{\pi(x)}x\Bigr)\longrightarrow0,
\]
as tails of convergent series (the last term accounts for the integer parts $\lfloor x/p\rfloor$ divided by $\lfloor x\rfloor$). By Volume~6, Theorem~9.1, applied to $\tilde f$, with $\tilde A=\sum_{y<p\le x}\frac{f^*(p)}p$ and $\tilde B^2=\sum_{y<p\le x}\frac{f^*(p)^2}p$,
\[
P_x\bigl(|\tilde f(n)-\tilde A|>\varepsilon\bigr)\le\frac{C_0\tilde B^2}{\varepsilon^2}\Bigl(1+O\Bigl(\frac1x\Bigr)\Bigr)\longrightarrow0,
\]
and $\tilde A\to0$, $\tilde B\to0$ as tails of the second and third series. So $P_x(|R(n)|>2\varepsilon)\to0$ for every $\varepsilon>0$.
\end{proof}

\begin{remark}[the converse]
Erd\H{o}s and Wintner also proved the converse: if $f$ has a limiting distribution, the three series converge. The proof needs a lower bound for characteristic functions of additive functions that we do not develop; see Tenenbaum's book, Chapter~III.4, or Elliott's \emph{Probabilistic Number Theory~I}, Chapter~5. Notice the asymmetry: the prime powers $p^k$ with $k\ge2$ play no role in the criterion, because $\sum_p\sum_{k\ge2}p^{-k}<\infty$.
\end{remark}

For $f(n)=\log\frac{\varphi(n)}n$ we have $f(p)=\log(1-\frac1p)\in[-\log2,0[$, so $f^*(p)=f(p)=-\frac1p+O(\frac1{p^2})$, and the three series converge (the first is empty). The same holds for $\log\frac{\sigma(n)}n$, with $f(p)=\log(1+\frac1p)$. So both have limiting distributions, and so do $\frac{\varphi(n)}n$ and $\frac{\sigma(n)}n$, by continuity of $\exp$. Figure~\ref{fig:phi} shows the first.

\begin{figure}[ht]
\centering\fig{ch01-phi}
\caption{The proportion of $n\le10^6$ with $\frac{\varphi(n)}n\le z$, against $z$. The limit law exists by Theorem~\ref{thm:erdoswintner} and is continuous by Corollary~\ref{cor:phicont}; the graph has no jumps, but its steep climbs at $z=\frac12$ (the even numbers not divisible by~$3$), $z=\frac13$ (multiples of $6$ prime to~$5$) and their neighbours show how singular it is. The mean is $\frac6{\pi^2}=0.6079$ (Exercise~\ref{ex:meanphi}).}
\label{fig:phi}
\end{figure}

\section{Continuity and purity of the limit law}

A sum of independent discrete variables is discrete when it has finitely many terms. With infinitely many it may be continuous: $\sum_k\varepsilon_k2^{-k}$, with independent fair $\varepsilon_k\in\{0,1\}$, is uniform on $[0,1]$. Paul L\'evy found in 1931 exactly when this happens.

\begin{think}
The limit law of $\frac{\varphi(n)}n$ is the law of $\prod_p(1-\frac1p)^{B_p}$, with independent $B_p$ equal to $1$ with probability about $\frac1p$. Does it have an atom, a value taken with positive probability? Compare with $\omega$-like functions where $f(p)=0$ for most primes.
\end{think}

The key is a lemma from the geometry of sums of signs.

\begin{lemma}[Littlewood--Offord, in Erd\H{o}s's form, 1945]\label{lem:lo}
Let $a_1,\dots,a_m$ be non-zero real numbers and $\varepsilon_1,\dots,\varepsilon_m$ independent with $P(\varepsilon_i=1)=P(\varepsilon_i=-1)=\frac12$. Then for every real $c$,
\[
P\Bigl(\sum_{i=1}^m\varepsilon_ia_i=c\Bigr)\le\binom m{\lfloor m/2\rfloor}2^{-m}\le\frac1{\sqrt m}.
\]
\end{lemma}
\begin{proof}
Replacing $(\varepsilon_i,a_i)$ by $(-\varepsilon_i,-a_i)$ does not change the law, so we may assume every $a_i>0$. A choice of signs is a subset $I=\{i:\varepsilon_i=1\}$ of $\{1,\dots,m\}$, and $\sum\varepsilon_ia_i=2\sum_{i\in I}a_i-\sum_ia_i$. If $I\subsetneq J$, the sum for $J$ exceeds the sum for $I$ by $2\sum_{J\setminus I}a_i>0$. So the sets $I$ giving the value $c$ form an \emph{antichain}: none contains another.

\emph{Sperner's bound.} An antichain $\mathcal F$ of subsets of an $m$-set has at most $\binom m{\lfloor m/2\rfloor}$ members. Indeed, consider the $m!$ maximal chains $\emptyset\subset\{\sigma(1)\}\subset\{\sigma(1),\sigma(2)\}\subset\cdots$, one for each permutation $\sigma$. A set $I$ of size $j$ lies on exactly $j!(m-j)!$ of them, and each chain meets the antichain at most once. Hence $\sum_{I\in\mathcal F}|I|!(m-|I|)!\le m!$, that is $\sum_{I\in\mathcal F}\binom m{|I|}^{-1}\le1$, and each $\binom m{|I|}\le\binom m{\lfloor m/2\rfloor}$.

\emph{The binomial bound.} For $m=2k$, $\bigl(\binom{2k}k4^{-k}\bigr)^2=\prod_{j=1}^k\bigl(\frac{2j-1}{2j}\bigr)^2\le\prod_{j=1}^k\frac{2j-1}{2j}\cdot\frac{2j}{2j+1}=\frac1{2k+1}$, so $\binom{2k}k4^{-k}\le\frac1{\sqrt{2k+1}}\le\frac1{\sqrt m}$. For $m=2k+1$, $\binom{2k+1}k2^{-2k-1}=\binom{2k}k4^{-k}\frac{2k+1}{2k+2}\le\frac1{\sqrt{2k+1}}=\frac1{\sqrt m}$.
\end{proof}

\begin{theorem}[L\'evy, 1931]\label{thm:levyatoms}
Let $X_1,X_2,\dots$ be independent discrete real random variables such that $S=\sum_jX_j$ converges almost surely, and let $q_j=\max_vP(X_j=v)$. Then the law of $S$ has an atom if and only if $\sum_j(1-q_j)<\infty$; otherwise it is continuous.
\end{theorem}

\begin{proof}
\emph{Convergent case.} Choose $v_j$ with $P(X_j=v_j)=q_j$. By the first Borel--Cantelli lemma, almost surely $X_j=v_j$ for all large $j$. Since $\sum X_j$ converges almost surely, $\sum_jv_j$ converges (its tail coincides with a tail of $\sum X_j$ on an event of positive probability). The events $E_N=\{X_j=v_j\ \forall j>N\}$ increase to an almost sure event, so $P(E_N)>0$ for some $N$. The finite sum $S_N=X_1+\dots+X_N$ is discrete, so $P(S_N=w)>0$ for some $w$. By independence,
\[
P\Bigl(S=w+\sum_{j>N}v_j\Bigr)\ge P(S_N=w)P(E_N)>0.
\]
\emph{Divergent case.} Let $S'$ be an independent copy of $S$, built from an independent copy $(X'_j)$. For every $a$,
\[
P(S=a)^2=P(S=a,S'=a)\le P(S-S'=0).
\]
Now $S-S'=\sum_jD_j$ with $D_j=X_j-X'_j$ independent and symmetric, and $P(D_j\ne0)=1-\sum_vP(X_j=v)^2\ge1-q_j$ since $\sum_vP(X_j=v)^2\le q_j\sum_vP(X_j=v)=q_j$. So $\sum_jP(D_j\ne0)=\infty$, and by the second Borel--Cantelli lemma infinitely many $D_j$ are non-zero, almost surely. Let $(\varepsilon_j)$ be independent fair signs, independent of everything else. By symmetry and independence, the sequence $(\varepsilon_jD_j)$ has the same law as $(D_j)$, so $P(\sum_jD_j=0)=P(\sum_j\varepsilon_jD_j=0)$.

Fix $m\ge1$ and $N$ so large that the event $F$ that at least $m$ of $D_1,\dots,D_N$ are non-zero has $P(F)\ge1-\frac1m$. Condition on all the $D_j$ and on the $\varepsilon_j$ with $D_j=0$ or $j>N$. On $F$, the sum $\sum_j\varepsilon_jD_j$ is then $c+\sum_{j\in J}\varepsilon_jD_j$ for a fixed $c$ and a set $J$ of at least $m$ indices with $D_j\ne0$, and the remaining signs are still independent and fair. Lemma~\ref{lem:lo} gives a conditional probability at most $\frac1{\sqrt m}$ that the sum is $0$. So $P(\sum D_j=0)\le\frac1{\sqrt m}+\frac1m$ for every $m$, hence $P(\sum D_j=0)=0$ and $S$ has no atom.
\end{proof}

\begin{corollary}\label{cor:ewcont}
Under the hypotheses of Theorem~\ref{thm:erdoswintner}, the limit law of $f$ is continuous if and only if $\sum_{f(p)\ne0}\frac1p=\infty$.
\end{corollary}
\begin{proof}
Let $X_p=f(p^{G_p})$. For $p\ge3$ every value other than $0$ is taken with probability at most $P(G_p\ge1)=\frac1p<1-\frac1p\le P(X_p=0)$, so $q_p=P(X_p=0)$ and $1-q_p=P(X_p\ne0)$. If $f(p)\ne0$ then $P(X_p\ne0)\ge P(G_p=1)=\frac{p-1}{p^2}\ge\frac1{2p}$; and always $P(X_p\ne0)\le\frac{\mathbf 1_{f(p)\ne0}}p+\frac1{p^2}$. So $\sum_p(1-q_p)$ and $\sum_{f(p)\ne0}\frac1p$ converge or diverge together, and Theorem~\ref{thm:levyatoms} applies.
\end{proof}

\begin{corollary}[Schoenberg 1928, Davenport 1933]\label{cor:phicont}
The functions $\frac{\varphi(n)}n$ and $\frac{\sigma(n)}n$ have continuous limiting distribution functions. In particular the set of abundant numbers, $\{n:\sigma(n)\ge2n\}$, has a natural density.
\end{corollary}
\begin{proof}
For both logarithms $f(p)\ne0$ for every $p$, and $\sum\frac1p=\infty$. Continuity of the limit law $F$ of $\log\frac{\sigma(n)}n$ at $\log2$ means that $P_x(\log\frac{\sigma(n)}n\ge\log2)\to1-F(\log2^-)=1-F(\log2)$.
\end{proof}

The argument shows more: the limit law of $\log\frac{\sigma(n)}n$ gives zero mass to every single point, so the density of $\{n:\sigma(n)\ge\alpha n\}$ exists for every $\alpha$ and depends continuously on $\alpha$. Figure~\ref{fig:sigma} shows the histogram up to $10^6$: $24.75\%$ of the integers are abundant there. The limit is $0.24761\ldots$, computed by Deléglise and then, to more digits, by Kobayashi and others, from the characteristic function of Theorem~\ref{thm:erdoswintner}.

\begin{figure}[ht]
\centering\fig{ch01-sigma}
\caption{Histogram of $\frac{\sigma(n)}n$ for $n\le10^6$ (density scale). The spike at $\frac{\sigma(n)}n=1^+$ comes from the primes and the integers with only large prime factors; abundant numbers are those right of the rust line. The largest value below $10^6$ is $4.509$, at $n=720720$.}
\label{fig:sigma}
\end{figure}

Continuity is not the end of the story. A probability law on $\R$ splits uniquely as $\mu=\mu_d+\mu_{sc}+\mu_{ac}$, a discrete part, a continuous part carried by a Lebesgue-null set, and an absolutely continuous part (Volume~7, Theorem~11.5; the singular part then splits into its atoms and the rest). Jessen and Wintner proved in 1935 that laws of convergent sums of independent discrete variables are always of one pure type.

\begin{theorem}[Jessen and Wintner, 1935]\label{thm:jessenwintner}
Under the hypotheses of Theorem~\ref{thm:levyatoms}, the law $\mu$ of $S$ is either purely discrete, or purely singular continuous, or absolutely continuous.
\end{theorem}
\begin{proof}
Let $G$ be the additive subgroup of $\R$ generated by all the values $v$ with $P(X_j=v)>0$, for all $j$; it is countable. For a Borel set $B$, the set $B+G=\bigcup_{g\in G}(B+g)$ is Borel. Almost surely every partial sum $S_N$ lies in $G$, so $S\in B+G$ if and only if $\sum_{j>N}X_j\in B+G$. This holds for every $N$, so up to a null set $\{S\in B+G\}$ is a tail event, and $\mu(B+G)\in\{0,1\}$ by Kolmogorov's zero--one law.

If $\mu_d\ne0$, let $D$ be the countable set of atoms; then $\mu(D+G)\ge\mu(D)>0$, so $\mu(D+G)=1$, and $D+G$ is countable: $\mu$ is purely discrete. If $\mu_d=0$ and $\mu_{sc}\ne0$, let $N$ be a Lebesgue-null set with $\mu_{sc}(N)=\mu_{sc}(\R)$. Then $N+G$ is null, a countable union of null sets, so $\mu_{ac}(N+G)=0$; and $\mu(N+G)\ge\mu_{sc}(N)>0$ forces $\mu(N+G)=1$, hence $\mu_{ac}(\R)=\mu_{ac}(\R)-\mu_{ac}(N+G)\le\mu(\R\setminus(N+G))=0$. So $\mu=\mu_{sc}$. Otherwise $\mu=\mu_{ac}$.
\end{proof}

\begin{remark}[which pure type?]
Erd\H{o}s showed in 1939 that the limit laws of $\frac{\varphi(n)}n$ and $\frac{\sigma(n)}n$ are purely singular: their distribution functions are continuous but have derivative $0$ almost everywhere. The steep climbs in Figure~\ref{fig:phi} are the visible trace of this. Deciding the type of the law of $\sum_p f(p^{G_p})$ in general is hard, and is related to the theory of Bernoulli convolutions; see Elliott, \emph{Probabilistic Number Theory~I}, Chapter~7.
\end{remark}

\begin{beauty}
Two unrelated-looking facts, that one integer in four is abundant and that $\frac{\varphi(n)}n$ has a continuous distribution, are instances of one probabilistic theorem about sums of independent coins. The arithmetic enters only through two counts: a prime $p$ divides a proportion $\frac1p$ of the integers, and the divisibilities by different small primes are independent up to an error that the sieve controls. Everything else is probability, and it is the probability of Kolmogorov, L\'evy and Sperner, a combinatorialist's lemma about subsets.
\end{beauty}


\section{The Selberg--Delange method}

\begin{center}\itshape
A square root of the zeta function sounds like nonsense.\\
Landau took one anyway, and the sums of two squares\\
came out thinned by exactly $\sqrt{\log x}$.
\end{center}

In 1908 Edmund Landau counted the integers that are sums of two squares. There are infinitely many of them, of course, and they have density zero; the question was how fast they thin out. Landau's answer was that their number up to $x$ is
\[
B(x)\sim b\,\frac x{\sqrt{\log x}},\qquad b=\frac1{\sqrt2}\prod_{p\equiv3\ (4)}\Bigl(1-\frac1{p^2}\Bigr)^{-1/2}=0.76422\ldots,
\]
and his proof contained a startling idea. The generating Dirichlet series of these integers behaves near $s=1$ like $\sqrt{\zeta(s)}$, which has no pole at $s=1$ but a \emph{branch point}, and Landau integrated around the branch cut.

Five years later, on 16 January 1913, a clerk at the Madras Port Trust sent a letter to G.~H.~Hardy in Cambridge. Among the dozens of formulas in it was the claim that $B(x)=b\int_2^x\frac{\dd t}{\sqrt{\log t}}+\theta(x)$ with $\theta(x)$ ``very small''. When Hardy asked what ``very small'' meant, Srinivasa Ramanujan wrote back: of order $\frac{\sqrt x}{\log x}$. That was false, as Daniel Shanks showed in 1964. But Hardy, relying on a computation by his student Geraldine Stanley that later turned out to be wrong, had written in his lectures on Ramanujan that Landau's approximation was the better one; Shanks showed that Ramanujan's is. Ramanujan was wrong in the right direction. Our own computer agrees: up to $10^7$ there are $1\,985\,459$ sums of two squares, Landau's formula gives $1\,903\,547$ and Ramanujan's gives $1\,969\,213$.

Forty years after Landau, in 1953 and 1954, the Indian mathematician L.~G.~Sathe published four long papers, more than a hundred and thirty pages in the \emph{Journal of the Indian Mathematical Society}, on a problem Hardy had raised: how many $n\le x$ have exactly $k$ prime factors, uniformly in $k$? Atle Selberg read them. In a five-page note printed in the same volume as Sathe's last two parts, he wrote that the results were very beautiful but the method complicated, and gave a new proof: study $\sum_{n\le x}z^{\omega(n)}$ for complex $z$, whose Dirichlet series is $\zeta(s)^z$ times a tame factor, and integrate around the branch cut of $\zeta(s)^z$ at $s=1$, exactly as Landau had done for $z=\frac12$. Hubert Delange turned the idea into a general theorem in 1959 and 1971, and the result is now called the \emph{Selberg--Delange method}.

Volume~10 (Chapter~11, equation~(11.1)) stated its conclusion for $\omega(n)$ and promised the proof here. The answer has a probabilistic face: the mean of $z^{\omega(n)}$ is the residue of a Poisson variable, corrected by an arithmetic factor and by $\frac1{\Gamma(z)}$, and that Gamma factor comes from the sharp cut-off $n\le x$ and from nothing else. Hankel's contour integral for $\frac1{\Gamma(z)}$ is the bridge.

\begin{reminder}[from Volumes~1, 4, 5 and~13]
\emph{Gamma} (Volume~4, Theorem~9.36). $\Gamma$ is meromorphic on $\C$ with poles at $0,-1,-2,\dots$, $\Gamma(z)\Gamma(1-z)=\frac\pi{\sin\pi z}$ for $z\notin\Z$, and $\frac1\Gamma$ is entire.
\emph{Zeros of $\zeta$} (Volume~5, Proposition~6.11 and Theorem~6.12). For $-1\le\sigma\le2$ and $|t|\ge2$, $\frac{\zeta'}\zeta(s)=\sum_{|\gamma-t|\le1}\frac1{s-\rho}+O(\log|t|)$, the sum having $O(\log|t|)$ terms; and $\zeta(\sigma+it)\ne0$ for $\sigma\ge1-\frac c{\log(|t|+2)}$, with an absolute $c>0$. Also $-\frac{\zeta'}\zeta(\sigma)\le\frac1{\sigma-1}+A$ for $1<\sigma\le2$.
\emph{Zeros of $L(s,\chi)$} (Volume~13, Lemma~2.2 and Theorem~2.3). For $\chi$ non-principal modulo $q$, $\mathcal L=\log(q(|t|+2))$ and $\frac78\le\sigma\le\frac{11}8$, $\frac{L'}L(s,\chi)=\sum_{|\rho-s_0|\le1/2}\frac1{s-\rho}+O(\mathcal L)$ with $s_0=\frac98+it$, the sum having $O(\mathcal L)$ terms; and $L(s,\chi)\ne0$ for $\sigma\ge1-\frac{c_1}{\mathcal L}$, except possibly for one real zero when $\chi$ is real.
\emph{Perron} (Volume~5, Lemma~7.1). For $y,c,T>0$ let $I(y,T)=\frac1{2\pi i}\int_{c-iT}^{c+iT}\frac{y^s}s\dd s$. Then $|I(y,T)-\mathbf 1_{y>1}|\le y^c\min(1,\frac1{\pi T|\log y|})$ for $y\ne1$, and $|I(1,T)-\frac12|\le\frac c{\pi T}$.
\emph{Two squares} (Volume~1, Theorem~7.3). $n$ is a sum of two squares if and only if every prime $q\equiv3\pmod4$ divides $n$ to an even power.
\end{reminder}

\section{Hankel's loop}

The whole method rests on one formula from the theory of the Gamma function. For $W>0$, let $\mathcal H_W$ be the path that comes from $-W$ along the lower edge of the negative real axis, goes once counterclockwise around the circle $|w|=1$, and returns to $-W$ along the upper edge; $\mathcal H=\mathcal H_\infty$ is \emph{Hankel's loop}. On it, $w^{-z}=e^{-z\log w}$ with the principal logarithm, $\arg w\in[-\pi,\pi]$, where $\arg w=-\pi$ on the lower edge and $\pi$ on the upper one.

\begin{lemma}[Hankel, 1864]\label{lem:hankel}
For every $z\in\C$,
\[
\frac1{2\pi i}\int_{\mathcal H}w^{-z}e^w\dd w=\frac1{\Gamma(z)}.
\]
Moreover, for $|z|\le R$ and $W\ge1$,
\[
\Bigl|\frac1{2\pi i}\int_{\mathcal H\setminus\mathcal H_W}w^{-z}e^w\dd w\Bigr|\le\frac{e^{\pi R}}\pi\int_W^\infty u^Re^{-u}\dd u.
\]
\end{lemma}
\begin{proof}
On the edges, $w=ue^{\pm i\pi}$ with $u\ge1$, and $|w^{-z}|=u^{-\operatorname{Re}z}e^{\pm\pi\operatorname{Im}z}\le e^{\pi R}u^R$ when $|z|\le R$, while $|e^w|=e^{-u}$. This gives the tail bound, and shows that the integral converges absolutely and uniformly for $|z|\le R$. The integrand is entire in $z$, so the integral is an entire function of $z$ (by Fubini and Morera's theorem, Volume~4, Chapter~9).

The circle $|w|=1$ can be replaced by $|w|=r$ for any $0<r\le1$: the difference is the integral over the boundary of a region where $w^{-z}e^w$ is holomorphic (the annulus slit along the negative axis), which vanishes by Cauchy's theorem. Let $\operatorname{Re}z<1$. On $|w|=r$ the integrand is at most $r^{-\operatorname{Re}z}e^{\pi|\operatorname{Im}z|}e^r$, and the circle has length $2\pi r$, so its contribution tends to $0$ as $r\to0$. On the lower edge, travelled from $-\infty$ to $0$, $w=ue^{-i\pi}$ and $w^{-z}=u^{-z}e^{i\pi z}$; on the upper edge, travelled from $0$ to $-\infty$, $w^{-z}=u^{-z}e^{-i\pi z}$. With $\dd w=-\dd u$,
\[
\int_{\mathcal H}w^{-z}e^w\dd w=e^{i\pi z}\int_0^\infty u^{-z}e^{-u}\dd u-e^{-i\pi z}\int_0^\infty u^{-z}e^{-u}\dd u=2i\sin(\pi z)\,\Gamma(1-z).
\]
So for $\operatorname{Re}z<1$ and $z\notin\Z$ the left side, divided by $2\pi i$, is $\frac{\sin\pi z}\pi\Gamma(1-z)=\frac1{\Gamma(z)}$ by the reflection formula. Both sides are entire, and they agree on an open set, so they agree everywhere.
\end{proof}

For $z=k$ a positive integer, $w^{-k}e^w$ is single-valued, the two edges cancel, and the lemma says that the residue of $e^ww^{-k}$ at $0$ is $\frac1{(k-1)!}$; for $z=0,-1,-2,\dots$ the integrand is entire and the integral is $0=\frac1{\Gamma(z)}$. Hankel's formula interpolates between these.

\begin{think}
Why should a sharp cut-off $n\le x$ produce $\frac1{\Gamma(z)}$? Compute $\frac1{2\pi i}\int_{c-i\infty}^{c+i\infty}(s-1)^{-z}\frac{x^s}s\dd s$ heuristically by the substitution $s=1+\frac w{\log x}$. What would the weight $e^{-n/x}$ instead of $\mathbf 1_{n\le x}$ give?
\end{think}

\section{Powers of zeta}

Let $c$ and $c_1$ be the absolute constants of Volume~5, Theorem~6.12, and Volume~13, Theorem~2.3. Put $c_2=\min(c,\frac{c_1}3)$ and $c_*=\min(\frac{c_2}2,\frac18)$, write $\tau=|t|+3$, and define the closed region and the open region
\[
\mathcal D=\Bigl\{s:\sigma\ge1-\frac{c_*}{\log\tau}\Bigr\}\subset U=\Bigl\{s:\sigma>1-\frac{2c_*}{\log\tau}\Bigr\}.
\]
Both lie in the half-plane $\sigma>\frac34$, because $\frac{2c_*}{\log3}\le\frac14\cdot\frac1{\log 3}<\frac14$. The region $U$ lies to the right of a curve, so it is simply connected.

Let $\chi_4$ be the non-principal character modulo $4$. We need three facts about $U$.
\begin{itemize}
\item $\zeta$ has no zero in $U$: there $\sigma>1-\frac{2c_*}{\log(|t|+3)}\ge1-\frac c{\log(|t|+2)}$.
\item $L(s,\chi_4)$ has no zero in $U$. Indeed $\log(4(|t|+2))\le3\log\tau$, so $U$ lies in the region of Volume~13, Theorem~2.3, with $q=4$; and the possible exceptional zero would be real, while $L(\sigma,\chi_4)=1-3^{-\sigma}+5^{-\sigma}-\cdots>0$ for real $\sigma>0$ (an alternating series with decreasing terms).
\item The function $Z(s)=(s-1)\zeta(s)$ is entire and has no zero in $U$, since $Z(1)=1$.
\end{itemize}
On the simply connected $U$, the non-vanishing holomorphic functions $Z$ and $L(\cdot,\chi_4)$ have holomorphic logarithms (Volume~4, Chapter~9), which we normalise to be real on $U\cap\R$. On the slit region $U^-=U\setminus\left]-\infty,1\right]$ put
\[
\log\zeta(s)=\log Z(s)-\log(s-1),\qquad \zeta(s)^z=e^{z\log\zeta(s)}=Z(s)^z(s-1)^{-z},
\]
with the principal $\log(s-1)$. This is a logarithm of $\zeta$, real on $\left]1,\infty\right[$, so for $\sigma>1$ it is the logarithm $\sum_n\frac{\Lambda(n)}{n^s\log n}=-\sum_p\log(1-p^{-s})$ of Volume~5, since two logarithms of the same function differ by a constant multiple of $2\pi i$ and these agree on the real axis. In particular, for $\sigma>1$, $\zeta(s)^z=\prod_p(1-p^{-s})^{-z}$ with principal powers, and $|\log\zeta(s)|\le\log\zeta(\sigma)$. Likewise $\log L(s,\chi_4)=-\sum_p\log(1-\chi_4(p)p^{-s})$ for $\sigma>1$, of modulus at most $\log\zeta(\sigma)$.

The growth of these logarithms is what the method needs. Write $f$ for either $\zeta$ or $L(\cdot,\chi_4)$.

\begin{lemma}\label{lem:logderivK}
There is an absolute constant $A_1$ such that $|\frac{f'}f(s)|\le A_1\log\tau$ for $s\in\mathcal D$ with $|t|\ge12$.
\end{lemma}
\begin{proof}
If $\sigma\ge2$, $|\frac{f'}f(s)|\le\sum_n\Lambda(n)n^{-2}$. Let $\sigma_1=1+\frac1{\log\tau}$ and $s_1=\sigma_1+it$; since $\log\tau\ge\log15>\frac83$, $\sigma_1<\frac{11}8$. For $\sigma\ge\sigma_1$, $|\frac{f'}f(s)|\le\sum_n\Lambda(n)n^{-\sigma}=-\frac{\zeta'}\zeta(\sigma)\le-\frac{\zeta'}\zeta(\sigma_1)\le\log\tau+A$.

Let $1-\frac{c_*}{\log\tau}\le\sigma\le\sigma_1$; then $\sigma\ge\frac78$. By the reminders, at $s$ and at $s_1$,
\[
\frac{f'}f(s)=\sum_{\rho\in\mathcal Z}\frac1{s-\rho}+O(\log\tau),\qquad\frac{f'}f(s_1)=\sum_{\rho\in\mathcal Z}\frac1{s_1-\rho}+O(\log\tau),
\]
where $\mathcal Z$ is a set of $O(\log\tau)$ zeros $\rho=\beta+i\gamma$ of $f$ with $|\gamma-t|\le1$ (the same set at $s$ and $s_1$, because it only depends on $t$). For such $\rho$, $\log(|\gamma|+2)\le\log\tau$ and, for $f=L(\cdot,\chi_4)$, $\log(4(|\gamma|+2))\le3\log\tau$; so the zero-free regions give $\beta\le1-\frac{c_2}{\log\tau}$.

\emph{Zeros seen from $s_1$.} Taking real parts at $s_1$, and using $|\frac{f'}f(s_1)|\le\log\tau+A$,
\[
\sum_{\rho\in\mathcal Z}\frac{\sigma_1-\beta}{|s_1-\rho|^2}=\operatorname{Re}\sum_{\rho\in\mathcal Z}\frac1{s_1-\rho}=O(\log\tau).
\]
Each $\sigma_1-\beta\ge\sigma_1-1=\frac1{\log\tau}$, so $\sum_{\rho\in\mathcal Z}|s_1-\rho|^{-2}=O(\log^2\tau)$.

\emph{From $s_1$ to $s$.} Since $\sigma\ge1-\frac{c_*}{\log\tau}\ge1-\frac{c_2}{2\log\tau}$, we have $\sigma-\beta\ge\frac{c_2}{2\log\tau}$, while $\sigma_1-\sigma\le\frac{1+c_*}{\log\tau}\le\frac2{\log\tau}$. So $\sigma_1-\beta\le(1+\frac4{c_2})(\sigma-\beta)$, and since $s$ and $s_1$ have the same imaginary part, $|s_1-\rho|\le(1+\frac4{c_2})|s-\rho|$. Subtracting the two expansions,
\[
\begin{aligned}\Bigl|\frac{f'}f(s)-\frac{f'}f(s_1)\Bigr|&\le\sum_{\rho\in\mathcal Z}\frac{|s_1-s|}{|s-\rho||s_1-\rho|}+O(\log\tau)\\&\le\frac2{\log\tau}\Bigl(1+\frac4{c_2}\Bigr)\sum_{\rho\in\mathcal Z}\frac1{|s_1-\rho|^2}+O(\log\tau)=O(\log\tau).\end{aligned}\qedhere
\]
\end{proof}

\begin{lemma}\label{lem:logzeta}
There is an absolute constant $A_2$ such that for $s\in\mathcal D$:
\begin{enumerate}[label=(\alph*)]
\item if $|t|\ge12$, then $|\log\zeta(s)|\le\log\log\tau+A_2$ and $|\log L(s,\chi_4)|\le\log\log\tau+A_2$;
\item if $|t|\le12$, then $|\log L(s,\chi_4)|\le A_2$, and $|\log Z(s)|\le A_2$ when moreover $\sigma\le2$.
\end{enumerate}
Consequently, for $|z|\le R$ and $|t|\ge12$, $|\zeta(s)^z|\le e^{RA_2}(\log\tau)^R$.
\end{lemma}
\begin{proof}
(a) For $\sigma\ge\sigma_1=1+\frac1{\log\tau}$, $|\log f(s)|\le\log\zeta(\sigma_1)\le\log\frac{\sigma_1}{\sigma_1-1}\le\log(2\log\tau)$, using $\zeta(\sigma)\le1+\int_1^\infty u^{-\sigma}\dd u=\frac\sigma{\sigma-1}$. For $1-\frac{c_*}{\log\tau}\le\sigma<\sigma_1$, $\log f(s)=\log f(s_1)-\int_\sigma^{\sigma_1}\frac{f'}f(u+it)\dd u$, and the integral is at most $\frac2{\log\tau}A_1\log\tau=2A_1$ by Lemma~\ref{lem:logderivK}. Take $A_2\ge\log2+2A_1$.

(b) The set $\{s\in\mathcal D:|t|\le12,\ \sigma\le2\}$ is compact and contained in $U$, where $\log Z$ and $\log L(\cdot,\chi_4)$ are continuous, hence bounded. For $\sigma\ge2$, $|\log L(s,\chi_4)|\le\log\zeta(2)$.

Finally $|\zeta(s)^z|=e^{\operatorname{Re}(z\log\zeta(s))}\le e^{R|\log\zeta(s)|}$.
\end{proof}

\section{The theorem}

For an integer $K\ge1$, let $\tau_K(n)$ be the number of ordered factorisations $n=d_1\cdots d_K$, so that $\sum_n\tau_K(n)n^{-s}=\zeta(s)^K$ for $\sigma>1$; it is multiplicative, and $\tau_K(p^a)$ is the number of $\mathbf u\in\N^K$ with $u_1+\dots+u_K=a$.

\begin{lemma}\label{lem:taupower}
For $K\ge1$, $n\ge1$ and $y\ge1$: $\tau_K(n)^2\le\tau_{K^2}(n)$, and $\sum_{n\le y}\tau_K(n)\le y(1+\log y)^K$.
\end{lemma}
\begin{proof}
Both sides of the first inequality are multiplicative, so let $n=p^a$. To a pair $(\mathbf u,\mathbf v)$ with $\sum u_i=\sum v_j=a$ we attach a $K\times K$ matrix of non-negative integers with row sums $\mathbf u$ and column sums $\mathbf v$; it exists by induction on $a$ (if $a\ge1$, pick $u_i\ge1$ and $v_j\ge1$, put a $1$ in position $(i,j)$, and apply the induction to $(\mathbf u-e_i,\mathbf v-e_j)$). Its row and column sums recover $(\mathbf u,\mathbf v)$, so the map is injective into the matrices with total $a$, which number $\tau_{K^2}(p^a)$.

For the second, $\sum_{n\le y}\tau_K(n)\le y\sum_{n\le y}\frac{\tau_K(n)}n\le y\bigl(\sum_{d\le y}\frac1d\bigr)^K$, since every factorisation $n=d_1\cdots d_K$ with $n\le y$ has all $d_i\le y$; and $\sum_{d\le y}\frac1d\le1+\log y$.
\end{proof}

\begin{theorem}[Selberg--Delange]\label{thm:sd}
Let $R\ge1$, $K\ge1$ and $1<q\le2$, and let $z\in\C$ with $|z|\le R$. Let $(a_n)$ be complex numbers with
\[
\sum_{n\le y}|a_n|^q\le Ky(1+\log y)^K\qquad(y\ge1),
\]
and $F(s)=\sum_na_nn^{-s}$ for $\sigma>1$. Suppose that $F(s)=\zeta(s)^zG(s)$ for $\sigma>1$, where $G$ is holomorphic on $U$ and $|G(s)|\le M(\log\tau)^R$ on $\mathcal D$. Then for $x\ge3$,
\[
\sum_{n\le x}a_n=x(\log x)^{z-1}\Bigl(\frac{G(1)}{\Gamma(z)}+O\Bigl(\frac{M+1}{\log x}\Bigr)\Bigr),
\]
where the implied constant depends only on $R$, $K$ and $q$.
\end{theorem}

Before the proof, a word on its shape. The sum is a Perron integral of $F(s)\frac{x^s}s$ along a vertical line. Pushing the line to the left of $1$ is impossible because of the branch cut of $\zeta(s)^z$, so the line wraps around the cut (Figure~\ref{fig:hankel}). Everything away from $s=1$ is small because $x^s$ is small there; the loop around the cut, rescaled by $\log x$, becomes Hankel's loop.

\begin{figure}[ht]
\centering\fig{ch02-hankel}
\caption{The contour of the proof. The Perron segment on $\sigma=c$ is replaced by the three sides of a rectangle, with a loop around the cut $[\sigma_T,1]$ of $\zeta(s)^z$. After the change of variables $s=1+\frac w{\log x}$ the loop becomes the truncated Hankel loop $\mathcal H_W$.}
\label{fig:hankel}
\end{figure}

\begin{proof}
Write $L=\log x$. All constants $C_1,C_2,\dots$ and all $O$'s depend only on $R$, $K$ and $q$. Put $\eta=\frac{q-1}{2q}$, $B=\frac{K+R+4}\eta$ and $T=L^B$. We first assume $x\ge x_0(R,K,q)$, where $x_0$ is so large that $T\le x$, $T\ge12$, $\log(T+3)\ge4$ and $\frac1L<\frac{c_*}{\log(T+3)}$; the small values of $x$ are treated at the end.

\emph{Step 1: Perron's formula.} Let $c=1+\frac1L$ and $A(y)=\sum_{n\le y}|a_n|$. Since $|a_n|\le1+|a_n|^q$, $A(y)\le2Ky(1+\log y)^K$, and by partial summation, with $y=e^v$,
\[
\sum_n\frac{|a_n|}{n^c}=c\int_1^\infty\frac{A(y)}{y^{c+1}}\dd y\le4K\int_0^\infty(1+v)^Ke^{-v/L}\dd v\le C_0L^{K+1}.
\]
So we may integrate term by term, and $\frac1{2\pi i}\int_{c-iT}^{c+iT}F(s)\frac{x^s}s\dd s=\sum_na_nI(\frac xn,T)$. The difference $E_1$ between this and $\sum_{n\le x}a_n$ is at most $\sum_n|a_n||\mathbf 1_{n\le x}-I(\frac xn,T)|$. Let $u=T^{-1/2}$. For $|\log\frac xn|\ge u$, Perron's lemma bounds the term by $|a_n|(\frac xn)^c\frac1{\pi Tu}$, and these terms add up to at most
\[
\frac{x^c}{\sqrt T}\sum_n\frac{|a_n|}{n^c}\le\frac{eC_0\,xL^{K+1}}{\sqrt T}.
\]
For $|\log\frac xn|<u$, the term is at most $2|a_n|$ (for $n\ne x$, $(\frac xn)^c\le e^{2u}<2$; for $n=x$, $|1-I(1,T)|\le\frac12+\frac c{\pi T}<1$). These $n$ lie in an interval of length $x(e^u-e^{-u})\le3xu$, so there are at most $3xu+1\le4xu$ of them (as $xu\ge\sqrt x\ge1$). By Hölder's inequality with exponents $\frac q{q-1}$ and $q$, and since $1+\log2x\le2L$,
\[
\sum_{|\log\frac xn|<u}|a_n|\le(4xu)^{1-\frac1q}\Bigl(\sum_{n\le2x}|a_n|^q\Bigr)^{\frac1q}\le(4xu)^{1-\frac1q}\bigl(2Kx(2L)^K\bigr)^{\frac1q}\le4K\,2^K\,xT^{-\eta}L^K.
\]
Altogether, as $\eta\le\frac14$, $|E_1|\le C_1xL^{K+1}T^{-\eta}=C_1xL^{-R-3}$.

\emph{Step 2: the contour.} Let $\delta=\frac{c_*}{\log(T+3)}$ and $\sigma_T=1-\delta$. For $|t|\le T$, $1-\frac{c_*}{\log\tau}\le\sigma_T$, so the rectangle $[\sigma_T,c]\times[-T,T]$ lies in $\mathcal D$, where $\zeta(s)^zG(s)\frac{x^s}s$ is holomorphic except on the cut $[\sigma_T,1]$; it coincides with $F(s)\frac{x^s}s$ on $\sigma=c$. By Cauchy's theorem, the integral over $[c-iT,c+iT]$ equals the sum of the integrals over the two horizontal sides $J_\pm$, the two vertical half-sides $V_\pm$ on $\sigma=\sigma_T$, and the loop $\Lambda$ around the cut made of its two edges and the circle $|s-1|=\frac1L$, all as in Figure~\ref{fig:hankel}. (Apply Cauchy's theorem to the rectangle minus a thin neighbourhood of the cut and let it shrink; the integrand is continuous up to the cut from each side.)

\emph{Step 3: the far part.} On $J_\pm$, $|t|=T\ge12$, so by Lemma~\ref{lem:logzeta} $|\zeta(s)^zG(s)|\le e^{RA_2}M(\log(T+3))^{2R}$, while $|\frac{x^s}s|\le\frac{ex}T$ and the length is at most $1$. So $|J_\pm|\le C_1Mx(\log(T+3))^{2R}T^{-1}$.

On $V_\pm$, $|x^s|=x^{\sigma_T}=xe^{-c_*L/\log(T+3)}$. Where $|t|\ge12$, the bound of $\zeta^zG$ is the same as on $J_\pm$. Where $|t|\le12$, $|\zeta(s)^z|=|Z(s)^z||(s-1)^{-z}|\le e^{RA_2}e^{\pi R}\max(|s-1|,|s-1|^{-1})^R\le C_2(\log(T+3))^R$, because $\delta\le|s-1|\le13$, and $|G(s)|\le M(\log 15)^R$. Since $|s|\ge\max(\frac12,|t|)$, $\int_{-T}^T\frac{\dd t}{|s|}\le4+2\log T$. So
\[
|V_+|+|V_-|\le C_3Mx(\log(T+3))^{2R+1}e^{-c_*L/\log(T+3)}.
\]
Now $\log(T+3)\le2B\log L$, so $e^{-c_*L/\log(T+3)}\le e^{-c_*L/(2B\log L)}$, which is $O(L^{-A})$ for every fixed $A$. Hence $|J_\pm|+|V_\pm|\le C_4MxL^{-R-3}$.

\emph{Step 4: the loop.} Write $\zeta(s)^zG(s)\frac1s=(s-1)^{-z}H(s)$ with $H(s)=Z(s)^z\frac{G(s)}s$. Let $r_0=\frac{c_*}2$. The disc $|s-1|\le r_0$ lies in $\mathcal D$ and in the compact set of Lemma~\ref{lem:logzeta}(b), so there $|H(s)|\le2e^{RA_2}(\log 4)^RM\le C_5M$; by Cauchy's estimate $|H'(s)|\le\frac{2C_5M}{r_0}$ on $|s-1|\le\frac{r_0}2$, so
\[
|H(s)-H(1)|\le\frac{2C_5M}{r_0}|s-1|\qquad(|s-1|\le\tfrac{r_0}2),\qquad H(1)=G(1),
\]
since $Z(1)=1$ and $\log Z(1)=0$. The loop $\Lambda$ lies in $|s-1|\le\delta\le\frac{r_0}2$, because $\log(T+3)\ge4$. Substitute $s=1+\frac wL$. The loop becomes $\mathcal H_W$ with $W=\delta L$, $x^s=xe^w$, $\dd s=\frac{\dd w}L$, and $(s-1)^{-z}=L^zw^{-z}$ (the principal arguments of $\frac wL$ and $w$ coincide). So
\[
\frac1{2\pi i}\int_\Lambda=xL^{z-1}\,\frac1{2\pi i}\int_{\mathcal H_W}w^{-z}e^wH\Bigl(1+\frac wL\Bigr)\dd w=xL^{z-1}\bigl(G(1)\,\mathcal I_1+\mathcal I_2\bigr),
\]
with
\[
\mathcal I_1=\frac1{2\pi i}\int_{\mathcal H_W}w^{-z}e^w\dd w,\qquad|\mathcal I_2|\le\frac{C_5M}{\pi r_0L}\int_{\mathcal H_W}|w||w^{-z}||e^w||\dd w|.
\]
On the circle $|w|=1$ the integrand of $\mathcal I_2$ is at most $e^{\pi R+1}$; on the edges it is at most $e^{\pi R}u^{R+1}e^{-u}$. So $|\mathcal I_2|\le C_6\frac ML$. By Lemma~\ref{lem:hankel}, $\mathcal I_1=\frac1{\Gamma(z)}+O(W^Re^{-W})$, and $W=\frac{c_*L}{\log(T+3)}\ge\frac{c_*L}{2B\log L}$, so $W^Re^{-W}=O(L^{-1})$. Also $|G(1)|\le M(\log3)^R\le2^RM$.

\emph{Step 5: conclusion for large $x$.} Collecting,
\[
\sum_{n\le x}a_n=xL^{z-1}\Bigl(\frac{G(1)}{\Gamma(z)}+O\Bigl(\frac ML\Bigr)\Bigr)+O\bigl((M+1)xL^{-R-3}\bigr),
\]
and $|xL^{z-1}|=xL^{\operatorname{Re}z-1}\ge xL^{-R-1}$, so the last error is $O((M+1)|xL^{z-1}|L^{-2})$.

\emph{Small $x$.} For $3\le x<x_0$, $L$ lies in the fixed interval $[\log3,\log x_0]$. The left side is at most $A(x)\le2Kx(1+L)^K$, and $|xL^{z-1}\frac{G(1)}{\Gamma(z)}|\le C_7Mx$ since $\frac1\Gamma$ is bounded on $|z|\le R$. So the difference between the two sides is $O((M+1)x)$, which is $O((M+1)|xL^{z-1}|L^{-1})$ because $|L^{z-1}|L^{-1}\ge L^{-R-2}$ is bounded below.
\end{proof}

\begin{remark}
Delange's theorem is sharper: under the same hypotheses there is an asymptotic expansion $\sum_{n\le x}a_n=xL^{z-1}\bigl(\sum_{j<J}\frac{\lambda_j(z)}{L^j}+O(L^{-J})\bigr)$, with $\lambda_j(z)=\frac1{\Gamma(z-j)}\frac1{j!}\frac{\dd^j}{\dd s^j}\bigl[(s-1)^z\zeta(s)^z\frac{G(s)}s\bigr]_{s=1}$, obtained by expanding $H$ to order $J$ in Step~4 (Tenenbaum, Chapter~II.5). The term $j=1$ explains the slow convergence in Figure~\ref{fig:sdratio}: for $z=3$ the ratio is still $1.45$ at $x=10^7$.
\end{remark}

\section{Prime factors: the promise of Volume~10}

\begin{lemma}\label{lem:eulerprod}
Let $R\ge1$ and $|z|\le R$. For $\sigma>1$, $\sum_nz^{\omega(n)}n^{-s}=\zeta(s)^zG_z(s)$, where
\[
G_z(s)=\prod_p\Bigl(1+\frac z{p^s-1}\Bigr)\Bigl(1-\frac1{p^s}\Bigr)^z
\]
converges absolutely and uniformly on $\sigma\ge\frac34$, is holomorphic on $\sigma>\frac34$, and satisfies $|G_z(s)|\le M(R)$ there. At $s=1$,
\[
G_z(1)=\prod_p\Bigl(1+\frac{z-1}p\Bigr)\Bigl(1-\frac1p\Bigr)^{z-1}.
\]
\end{lemma}
\begin{proof}
Let $K$ be the least integer $\ge R$. Since $|z^{\omega(n)}|\le K^{\omega(n)}\le\tau_K(n)$, the series converges absolutely for $\sigma>1$ and, $z^{\omega(n)}$ being multiplicative, it has the Euler product $\prod_p(1+z\sum_{k\ge1}p^{-ks})=\prod_p\phi(p^{-s})$, where $\phi(u)=1+\frac{zu}{1-u}=(1-u+zu)(1-u)^{-1}$. Multiplying by $\prod_p(1-p^{-s})^z=\zeta(s)^{-z}$ gives the identity, with factors $g_p(s)=\psi(p^{-s})$ where
\[
\psi(u)=(1+(z-1)u)(1-u)^{z-1},\qquad|u|<1.
\]
Now $\psi(0)=1$ and $\psi'(0)=(z-1)-(z-1)=0$, so $\frac{\psi(u)-1}{u^2}$ is holomorphic on $|u|<1$. On $|u|=0.8$, $|1-u|\ge0.2$ and $|\log(1-u)|\le\log5+\frac\pi2$, so $|\psi(u)-1|\le1+1.8(R+1)e^{(R+1)(\log5+\pi/2)}=:C_R$; by the maximum principle, $|\psi(u)-1|\le\frac{C_R}{0.64}|u|^2$ for $|u|\le0.8$. For $\sigma\ge\frac34$, $|p^{-s}|\le2^{-3/4}<0.6$, so $|g_p(s)-1|\le2C_Rp^{-3/2}$. Since $\sum_pp^{-3/2}<2$, the product converges absolutely and uniformly, its limit is holomorphic, and $|G_z(s)|\le\prod_p(1+2C_Rp^{-3/2})\le e^{4C_R}$. Finally $(1+\frac z{p-1})(1-\frac1p)^z=\frac{p-1+z}p(1-\frac1p)^{z-1}$.
\end{proof}

\begin{corollary}[Selberg, 1954]\label{cor:omegasd}
Uniformly for $|z|\le R$ and $x\ge3$,
\[
\frac1x\sum_{n\le x}z^{\omega(n)}=(\log x)^{z-1}\Bigl(\frac{F(z)}{\Gamma(z)}+O_R\Bigl(\frac1{\log x}\Bigr)\Bigr),\qquad F(z)=\prod_p\Bigl(1+\frac{z-1}p\Bigr)\Bigl(1-\frac1p\Bigr)^{z-1}.
\]
This is equation~(11.1) of Volume~10.
\end{corollary}
\begin{proof}
Apply Theorem~\ref{thm:sd} with $a_n=z^{\omega(n)}$, $q=2$ and $G=G_z$. With $K_0$ the least integer $\ge R$, $|a_n|^2\le\tau_{K_0}(n)^2\le\tau_{K_0^2}(n)$, so Lemma~\ref{lem:taupower} gives the hypothesis on $(a_n)$ with $K=K_0^2$; and since $U\subset\{\sigma>\frac34\}$, Lemma~\ref{lem:eulerprod} gives the hypothesis on $G$ with $M=M(R)$.
\end{proof}

In the probabilistic language of Volume~10: for $n$ uniform in $[1,x]$, $\mathbb E_xz^{\omega(n)}=e^{(z-1)\log\log x}\bigl(\frac{F(z)}{\Gamma(z)}+o(1)\bigr)$, the generating function of a Poisson variable of parameter $\log\log x$, times a residual factor. The next section extracts from it the sharp form of the Erd\H{o}s--Kac theorem and the Sathe--Selberg formula for $\#\{n\le x:\omega(n)=k\}$.

\begin{figure}[ht]
\centering\fig{ch02-ratio}
\caption{The ratio of $\sum_{n\le x}z^{\omega(n)}$ to the main term $x(\log x)^{z-1}\frac{F(z)}{\Gamma(z)}$ of Corollary~\ref{cor:omegasd}, for $z=\frac12$, $2$ and $3$, with $x$ from $10^3$ to $10^7$. The ratios tend to $1$, as slowly as $1+\frac{\text{const}}{\log x}$ predicts: at $10^7$ they are $1.034$, $1.080$ and $1.454$. For $z=2$, $F(2)=\prod_p(1-\frac1{p^2})=\frac6{\pi^2}$.}
\label{fig:sdratio}
\end{figure}

The same method counts prime factors with multiplicity, as long as $|z|<2$.

\begin{corollary}\label{cor:bigomegasd}
Uniformly for $|z|\le\frac32$ and $x\ge3$,
\[
\sum_{n\le x}z^{\Omega(n)}=x(\log x)^{z-1}\Bigl(\frac{F_\Omega(z)}{\Gamma(z)}+O\Bigl(\frac1{\log x}\Bigr)\Bigr),\qquad F_\Omega(z)=\prod_p\Bigl(1-\frac1p\Bigr)^z\Bigl(1-\frac zp\Bigr)^{-1}.
\]
\end{corollary}
\begin{proof}
\emph{The coefficients.} Here $|z^{\Omega(n)}|$ can be as large as $(\frac32)^{\Omega(n)}$, which is not bounded by any $\tau_K(n)$ (take $n=2^a$); this is why Theorem~\ref{thm:sd} allows $q<2$. Take $q=\frac32$ and $w=(\frac32)^{3/2}<1.84$. Every $n\le y$ factors over the primes $p\le y$, and $\sum_kw^kp^{-k}=(1-\frac wp)^{-1}$, so
\[
\sum_{n\le y}|z|^{\frac32\Omega(n)}\le y\sum_{n\le y}\frac{w^{\Omega(n)}}n\le y\prod_{p\le y}\Bigl(1-\frac wp\Bigr)^{-1}\le C\,y(1+\log y)^2,
\]
because $-\log(1-v)\le v+7v^2$ for $0\le v\le0.92$ and $\sum_{p\le y}\frac1p\le\log(1+\log y)+C$ (Volume~5, Theorem~1.5). This is the hypothesis of Theorem~\ref{thm:sd} with $K=\max(C,2)$.

\emph{The Euler factors.} Since $\Omega$ is completely additive, $\sum_nz^{\Omega(n)}n^{-s}=\prod_p(1-zp^{-s})^{-1}$ for $\sigma>1$, and multiplying by $\zeta(s)^{-z}=\prod_p(1-p^{-s})^z$ gives $G(s)=\prod_p\psi(p^{-s})$ with $\psi(u)=(1-u)^z(1-zu)^{-1}$. For $|u|<\frac23$, $|zu|<1$, so $\psi$ is holomorphic there, with $\psi(0)=1$ and $\psi'(0)=-z+z=0$. On $|u|=0.62$, $|1-zu|\ge0.07$ and $|(1-u)^z|\le e^{\frac32(\log3+\frac\pi2)}$, so $|\psi-1|$ is bounded by an absolute constant, and by the maximum principle $|\psi(u)-1|\le C'|u|^2$ for $|u|\le0.62$. For $\sigma\ge\frac34$, $|p^{-s}|\le2^{-3/4}<0.6$, and the proof of Lemma~\ref{lem:eulerprod} applies word for word: $G$ is holomorphic and bounded on $\sigma>\frac34$, and $G(1)=F_\Omega(z)$. Theorem~\ref{thm:sd} with $R=\frac32$ concludes.
\end{proof}

\section{Landau's sums of two squares}

\begin{corollary}[Landau, 1908]\label{cor:landau}
The number $B(x)$ of integers $n\le x$ that are sums of two squares satisfies
\[
B(x)=\frac{b\,x}{\sqrt{\log x}}\Bigl(1+O\Bigl(\frac1{\log x}\Bigr)\Bigr),\qquad b=\frac1{\sqrt2}\prod_{p\equiv3\ (4)}\Bigl(1-\frac1{p^2}\Bigr)^{-1/2}=0.764223\ldots
\]
\end{corollary}
\begin{proof}
Let $a_n=1$ if $n$ is a sum of two squares and $0$ otherwise. By Volume~1, Theorem~7.3, $a_n$ is multiplicative, and for $\sigma>1$
\[
F(s)=\sum_na_nn^{-s}=(1-2^{-s})^{-1}\prod_{p\equiv1\ (4)}(1-p^{-s})^{-1}\prod_{p\equiv3\ (4)}(1-p^{-2s})^{-1}.
\]
Let $G(s)=L(s,\chi_4)^{1/2}P(s)$ with $L^{1/2}=e^{\frac12\log L}$ and
\[
P(s)=(1-2^{-s})^{-1/2}\prod_{p\equiv3\ (4)}(1-p^{-2s})^{-1/2},
\]
principal powers; $P$ is holomorphic and $|P|\le(1-2^{-3/4})^{-1/2}\zeta(\frac32)^{1/2}$ on $\sigma>\frac34$. For $\sigma>1$, compare logarithms prime by prime in $\log F$ and $\frac12\log\zeta+\log G$: at $p=2$ both are $-\log(1-2^{-s})$; at $p\equiv1\pmod4$ both are $-\log(1-p^{-s})$; at $p\equiv3\pmod4$ the right side is $-\frac12\log(1-p^{-s})-\frac12\log(1+p^{-s})-\frac12\log(1-p^{-2s})=-\log(1-p^{-2s})$. So $F=\zeta^{1/2}G$.

Take $z=\frac12$, $R=K=1$ and $q=2$ (since $|a_n|\le1$). By Lemma~\ref{lem:logzeta}, $|L(s,\chi_4)^{1/2}|\le e^{A_2/2}(\log\tau)^{1/2}$ on $\mathcal D$ (for $|t|\le12$ use part~(b)), so $|G|\le M\log\tau$ for an absolute $M$. Theorem~\ref{thm:sd} gives $B(x)=x(\log x)^{-1/2}(\frac{G(1)}{\Gamma(1/2)}+O(\frac1{\log x}))$. Finally $L(1,\chi_4)=\frac\pi4$ (Leibniz; Volume~5, Example~3.11) and $\Gamma(\frac12)=\sqrt\pi$, so
\[
\frac{G(1)}{\sqrt\pi}=\frac1{\sqrt\pi}\Bigl(\frac\pi4\Bigr)^{1/2}\sqrt2\prod_{p\equiv3\ (4)}\Bigl(1-\frac1{p^2}\Bigr)^{-1/2}=b.\qedhere
\]
\end{proof}

\begin{beauty}
The constant $b$ holds three pieces of the integers in one formula. The $\frac1{\sqrt\pi}$ from $\Gamma(\frac12)$ is the cut-off, the same as for the cycles of a random permutation. The $\sqrt{\pi/4}$ is Leibniz's series, the value at $1$ of the $L$-function that tells apart the primes $1$ and $3$ modulo~$4$. And the product over $p\equiv3\pmod4$ counts the price of those primes having to come in pairs.
\end{beauty}

Ramanujan's version replaces $\frac x{\sqrt{\log x}}$ by $\int_2^x\frac{\dd t}{\sqrt{\log t}}=\frac x{\sqrt{\log x}}(1+\frac1{2\log x}+\cdots)$, which builds in a second-order term of Delange's expansion. Shanks computed the true second-order term. It is not the one the integral produces, so Ramanujan's claim about $\theta(x)$ is false; but it is closer to the integral's than to Landau's, which has none, so the integral is the better approximation for all large $x$. Exercise~\ref{ex:onlythree} counts the opposite set: integers all of whose prime factors are $3$ modulo~$4$.

\section{Exercises}

\begin{exercise}[$\star$]\label{ex:hankelint}
Check Lemma~\ref{lem:hankel} directly for $z=1,2,3$ by the residue theorem, and for $z=0,-1$ by Cauchy's theorem. Use it to compute $\frac1{2\pi i}\int_{\mathcal H}w^{-1/2}e^w\dd w$.
\end{exercise}

\begin{exercise}[$\star$]\label{ex:twoomega}
Deduce from Corollary~\ref{cor:omegasd} that $\sum_{n\le x}2^{\omega(n)}=\frac6{\pi^2}x\log x+O(x)$. Then prove it without complex analysis, from $2^{\omega(n)}=\sum_{d\mid n}\mu(d)^2$.
\end{exercise}

\begin{exercise}[$\star\star$]\label{ex:minusone}
Show that $\sum_{n\le x}(-1)^{\omega(n)}=O(\frac x{(\log x)^3})$. Which value of the Gamma function is responsible? Numerically, one finds $\frac1x\sum_{n\le x}(-1)^{\omega(n)}=-0.000165$ at $x=10^7$.
\end{exercise}

\begin{exercise}[$\star\star$]\label{ex:bigomegasd}
Show that $\sum_{n\le x}2^{\Omega(n)}\ge c\,x(\log x)^2$ for $x\ge2$ and some $c>0$, so that Corollary~\ref{cor:bigomegasd} cannot extend to $z=2$. Where does its proof break down?
\end{exercise}

\begin{exercise}[$\star\star$]\label{ex:onlythree}
Show that the number of $n\le x$ all of whose prime factors are $\equiv3\pmod4$ is $\frac{c_3x}{\sqrt{\log x}}(1+O(\frac1{\log x}))$, with $c_3=\frac{\sqrt2}\pi\prod_{p\equiv3\ (4)}(1-\frac1{p^2})^{-1/2}=0.48652\ldots$
\end{exercise}

\begin{exercise}[$\star\star\star$]\label{ex:landaupik}
(Landau, 1900.) Let $\pi_k(x)=\#\{n\le x:\omega(n)=k\}$. Show, for each fixed $k\ge1$, that
\[
\pi_k(x)=\frac x{\log x}\,\frac{(\log\log x)^{k-1}}{(k-1)!}\Bigl(1+O_k\Bigl(\frac1{\log\log x}\Bigr)\Bigr).
\]
\end{exercise}

\section{Hints}

\begin{hint}{ex:hankelint}
For $z=k\ge1$ the two edges cancel, since $w^{-k}$ is single-valued. For $z=\frac12$, $\Gamma(\frac12)=\sqrt\pi$.
\end{hint}
\begin{hint}{ex:twoomega}
$F(2)=\frac6{\pi^2}$ and $\Gamma(2)=1$. For the elementary proof, swap sums and use $\sum_{d\le y}\mu(d)^2=\frac6{\pi^2}y+O(\sqrt y)$.
\end{hint}
\begin{hint}{ex:minusone}
$\frac1{\Gamma(-1)}=0$.
\end{hint}
\begin{hint}{ex:bigomegasd}
Write $n=2^km$ with $m$ odd, and use $2^{\Omega(m)}\ge2^{\omega(m)}$ and Theorem~\ref{thm:sd} for the odd $m$.
\end{hint}
\begin{hint}{ex:onlythree}
Find $G$ with $F=\zeta^{1/2}G$; this time $L(s,\chi_4)$ appears with the exponent $-\frac12$.
\end{hint}
\begin{hint}{ex:landaupik}
Write
\[
\pi_k(x)=\frac1{2\pi i}\oint_{|z|=1/2}z^{-k-1}\sum_{n\le x}z^{\omega(n)}\dd z;
\]
the function $\Phi(z)=\frac{F(z)}{\Gamma(z)}$ is entire with $\Phi(z)=z+O(z^2)$.
\end{hint}

\section{Solutions}

\begin{solution}{ex:hankelint}
\idea Where $w^{-z}$ is single-valued, the loop can be closed into a circle and Hankel's formula becomes a residue computation.
\writeup For an integer $z$, $w^{-z}$ is single-valued, so the integrals along the two edges are equal and opposite and $\frac1{2\pi i}\int_{\mathcal H}w^{-z}e^w\dd w=\frac1{2\pi i}\oint_{|w|=1}w^{-z}e^w\dd w$. For $z=k\ge1$ this is the coefficient of $w^{k-1}$ in $e^w$, namely $\frac1{(k-1)!}=\frac1{\Gamma(k)}$: $1,1,\frac12$ for $k=1,2,3$. For $z=0$ or $-1$ the integrand $e^w$ or $we^w$ is entire and the integral vanishes, as does $\frac1\Gamma$ at $0$ and $-1$. For $z=\frac12$ the lemma gives $\frac1{\Gamma(1/2)}=\frac1{\sqrt\pi}$.
\end{solution}

\begin{solution}{ex:twoomega}
\idea Both proofs see the same constant: $\frac6{\pi^2}$ is the density of squarefree numbers, and $2^{\omega(n)}$ counts the squarefree divisors of $n$.
\writeup With $z=2$, $F(2)=\prod_p(1+\frac1p)(1-\frac1p)=\prod_p(1-\frac1{p^2})=\frac6{\pi^2}$ and $\Gamma(2)=1$, so Corollary~\ref{cor:omegasd} gives $x\log x(\frac6{\pi^2}+O(\frac1{\log x}))=\frac6{\pi^2}x\log x+O(x)$. Elementarily, $2^{\omega(n)}$ is the number of squarefree $d\mid n$, so $\sum_{n\le x}2^{\omega(n)}=\sum_{d\le x}\mu(d)^2\lfloor\frac xd\rfloor=x\sum_{d\le x}\frac{\mu(d)^2}d+O(x)$. By partial summation from $Q(y)=\sum_{d\le y}\mu(d)^2=\frac6{\pi^2}y+O(\sqrt y)$, $\sum_{d\le x}\frac{\mu(d)^2}d=\frac{Q(x)}x+\int_1^x\frac{Q(y)}{y^2}\dd y=\frac6{\pi^2}\log x+O(1)$.
\end{solution}

\begin{solution}{ex:minusone}
\idea At $z=-1$ the main term of the Selberg--Delange formula vanishes because $\frac1\Gamma$ has a zero at $-1$; only the error term remains.
\writeup Corollary~\ref{cor:omegasd} with $z=-1$, $R=1$: $\sum_{n\le x}(-1)^{\omega(n)}=x(\log x)^{-2}(\frac{F(-1)}{\Gamma(-1)}+O(\frac1{\log x}))$, and $\frac1{\Gamma(-1)}=0$, so the sum is $O(\frac x{(\log x)^3})$. In the language of the proof: $\zeta(s)^{-1}$ has no branch cut, the loop closes up, and the integrand $\zeta(s)^{-1}G(s)\frac{x^s}s$ is holomorphic at $s=1$. Delange's expansion of the Remark says more: every term contains $\frac1{\Gamma(-1-j)}=0$, so the sum is $O(\frac x{(\log x)^A})$ for every $A$. To do better one must use the zero-free region itself, as for the prime number theorem.
\end{solution}

\begin{solution}{ex:bigomegasd}
\idea The prime $2$ alone already does the damage: the powers of $2$ carry a weight $2^k$ that compensates exactly for their rarity.
\writeup For odd $m$, $a_m=2^{\omega(m)}\mathbf 1_{m\text{ odd}}$ has the Dirichlet series $\prod_{p>2}(1+\frac2{p^s-1})=\zeta(s)^2G_2(s)\frac{2^s-1}{2^s+1}$, with $G_2$ as in Lemma~\ref{lem:eulerprod}; the last factor is holomorphic and bounded on $\sigma>\frac34$ (there $|2^s+1|\ge2^{3/4}-1$) and equals $\frac13$ at $s=1$. Theorem~\ref{thm:sd} ($z=2$, $q=2$) gives $\sum_{m\le y,\,m\text{ odd}}2^{\omega(m)}\sim\frac{F(2)}3y\log y=\frac2{\pi^2}y\log y$, so this sum is at least $c\,y\log2y$ for all $y\ge1$ (for small $y$ the term $m=1$ suffices). Now $2^{\Omega(2^km)}=2^k2^{\Omega(m)}\ge2^k2^{\omega(m)}$, and with $K=\lfloor\log_2x\rfloor$,
\[
\sum_{n\le x}2^{\Omega(n)}\ge\sum_{k=0}^K2^k\,c\,\frac x{2^k}\log\frac{2x}{2^k}\ge c\,x\sum_{k=0}^K(K+1-k)\log2\ge\frac c{2\log2}\,x(\log x)^2.
\]
The proof of Corollary~\ref{cor:bigomegasd} breaks twice at $z=2$: the Euler factor at $p=2$ becomes $(1-2^{1-s})^{-1}$, with a pole at $s=1$, so $G$ is not holomorphic there; and no $q>1$ makes $\sum_{n\le y}2^{q\Omega(n)}$ of size $y(\log y)^K$, since the powers of $2$ alone give $y^q$.
\end{solution}

\begin{solution}{ex:onlythree}
\idea Write the Euler product, compare it with $\zeta(s)/L(s,\chi_4)$, and take a square root.
\writeup The series is $F(s)=\prod_{p\equiv3\ (4)}(1-p^{-s})^{-1}$. Prime by prime, $\frac{\zeta(s)}{L(s,\chi_4)}$ has the factor $(1-2^{-s})^{-1}$ at $2$, $1$ at $p\equiv1$, and $\frac{1+p^{-s}}{1-p^{-s}}$ at $p\equiv3$. So $F^2=(1-2^{-s})\frac\zeta L\prod_{p\equiv3}(1-p^{-2s})^{-1}$, and comparing logarithms as in Corollary~\ref{cor:landau}, $F=\zeta^{1/2}G$ with $G=(1-2^{-s})^{1/2}L(s,\chi_4)^{-1/2}\prod_{p\equiv3}(1-p^{-2s})^{-1/2}$. Lemma~\ref{lem:logzeta} bounds $|L^{-1/2}|$ by $e^{A_2/2}(\log\tau)^{1/2}$, so Theorem~\ref{thm:sd} applies with $z=\frac12$, and the constant is $\frac{G(1)}{\sqrt\pi}=\frac1{\sqrt\pi}\cdot\frac1{\sqrt2}\cdot\frac2{\sqrt\pi}\prod_{p\equiv3}(1-p^{-2})^{-1/2}=\frac{\sqrt2}\pi\prod_{p\equiv3}(1-p^{-2})^{-1/2}$. Numerically, one finds the ratio of the count to $\frac{c_3x}{\sqrt{\log x}}$ to be $1.026$ at $x=10^7$.
\end{solution}

\begin{solution}{ex:landaupik}
\idea The Selberg--Delange formula holds uniformly on a circle around $z=0$, so Cauchy's formula extracts the coefficient of $z^k$; the main term is a Taylor coefficient of $e^{z\log\log x}\Phi(z)$.
\writeup Let $\lambda=\log\log x$ and $L=\log x$. Since $\sum_{n\le x}z^{\omega(n)}=\sum_k\pi_k(x)z^k$ is a polynomial,
\[
\pi_k(x)=\frac1{2\pi i}\oint_{|z|=1/2}\frac{\sum_{n\le x}z^{\omega(n)}}{z^{k+1}}\dd z=\frac xL\,\frac1{2\pi i}\oint_{|z|=1/2}\frac{e^{\lambda z}\Phi(z)}{z^{k+1}}\dd z+E,
\]
with $\Phi=\frac F\Gamma$ and, by Corollary~\ref{cor:omegasd} with $R=1$, $|E|\le\frac1{2\pi}\pi\,2^{k+1}\max_{|z|=1/2}|xL^{z-1}|\cdot O(\frac1L)=O_k(xL^{-3/2})$. The main integral is the coefficient of $z^k$ in $e^{\lambda z}\Phi(z)$. The function $\Phi$ is entire, with $\Phi(z)=\sum_{j\ge1}\phi_jz^j$, because $\frac1{\Gamma(z)}=z+O(z^2)$ and $F(0)=\prod_p(1-\frac1p)(1-\frac1p)^{-1}=1$; so $\phi_1=1$. The coefficient is $\sum_{j=1}^k\phi_j\frac{\lambda^{k-j}}{(k-j)!}=\frac{\lambda^{k-1}}{(k-1)!}(1+O_k(\frac1\lambda))$. Since $xL^{-3/2}=\frac xL\cdot L^{-1/2}$ is much smaller than $\frac xL\frac{\lambda^{k-2}}{(k-1)!}$, the claim follows. For $k=1$ this is the prime number theorem with $\pi_1(x)=\pi(x)+\#\{p^a\le x,a\ge2\}$, in the weak form $\frac x{\log x}(1+O(\frac1{\log\log x}))$.
\end{solution}

\section*{Sources and further reading}

\begin{itemize}
\item \emph{Landau and Ramanujan.} E.~Landau, ``Über die Einteilung der positiven ganzen Zahlen in vier Klassen nach der Mindestzahl der zu ihrer additiven Zusammensetzung erforderlichen Quadrate'', \emph{Arch.\ Math.\ Phys.}\ (3) 13 (1908), 305--312. B.~C.~Berndt and R.~A.~Rankin, \emph{Ramanujan: Letters and Commentary}, AMS/LMS, 1995, for the letter of 16 January 1913. G.~H.~Hardy, \emph{Ramanujan: Twelve Lectures on Subjects Suggested by His Life and Work}, Cambridge, 1940, pp.~9 and~63. D.~Shanks, ``The second-order term in the asymptotic expansion of $B(x)$'', \emph{Math.\ Comp.}\ 18 (1964), 75--86. P.~Moree, ``Counting numbers in multiplicative sets: Landau versus Ramanujan'', \emph{Mathematics Newsletter} 21 (2011), 73--81 (arXiv:1110.0708), for the story of Ramanujan's reply, of Stanley's computation and of Shanks's verdict; P.~Moree and J.~Cazaran, ``On a claim of Ramanujan in his first letter to Hardy'', \emph{Expo.\ Math.}\ 17 (1999), 289--312.
\item \emph{Sathe and Selberg.} L.~G.~Sathe, ``On a problem of Hardy on the distribution of integers having a given number of prime factors'', I--II, \emph{J.\ Indian Math.\ Soc.}\ 17 (1953), 63--141; III--IV, 18 (1954), 27--81. A.~Selberg, ``Note on a paper by L.~G.~Sathe'', \emph{J.\ Indian Math.\ Soc.}\ 18 (1954), 83--87; its opening paragraph calls Sathe's results very beautiful and his method complicated and involved.
\item \emph{Delange.} H.~Delange, ``Sur des formules dues à Atle Selberg'', \emph{Bull.\ Sci.\ Math.}\ (2) 83 (1959), 101--111; ``Sur des formules de Atle Selberg'', \emph{Acta Arith.}\ 19 (1971), 105--146.
\item \emph{Landau 1900.} E.~Landau, ``Sur quelques problèmes relatifs à la distribution des nombres premiers'', \emph{Bull.\ Soc.\ Math.\ France} 28 (1900), 25--38, for $\pi_k(x)$.
\item \emph{Hankel.} H.~Hankel, ``Die Euler'schen Integrale bei unbeschränkter Variabilität des Arguments'', \emph{Z.\ Math.\ Phys.}\ 9 (1864), 1--21.
\item \emph{Textbooks.} G.~Tenenbaum, \emph{Introduction to Analytic and Probabilistic Number Theory}, 3rd ed., AMS, 2015, Chapters~II.5--II.6, whose formulation of the method we follow. H.~L.~Montgomery and R.~C.~Vaughan, \emph{Multiplicative Number Theory~I}, Cambridge, 2007, Chapter~7. R.~de~la~Bretèche and G.~Tenenbaum, ``Remarks on the Selberg--Delange method'', \emph{Acta Arith.}\ 200 (2021), 349--369 (arXiv:2010.12929).
\end{itemize}


\section{The Erd\H{o}s--Kac theorem, as a corollary}

The Erd\H os--Kac theorem is a corollary of the Kubilius model, and it is already a theorem of Volume~6. One proof is enough to cite. This section says what the corollary is, in two pages, and what each proof makes visible. It does not found the theorem a second time. The Sathe--Selberg formula, which sharpens the count of integers with $\omega(n)=k$, is Cauchy's integral of the Selberg--Delange mean already proved in this chapter; it is not rewritten here.

\begin{theorem}[Erd\H os--Kac, as a corollary]\label{thm:erdoskac-cor}
For every real $a$, the proportion of $n\le x$ with $(\omega(n)-\log\log x)/\sqrt{\log\log x}\le a$ tends to $\Phi(a)$.
\end{theorem}
\begin{proof}
Let $L=\log\log x$ and $y=\exp(L^{1/2})$, so $\log\log y=\tfrac12 L\to\infty$ and $u=\log x/\log y=L^{1/2}\to\infty$. Theorem~\ref{thm:kubilius} gives total variation $o(1)$ between the law of $(\nu_p(n))_{p\le y}$ and the independent geometrics $G_p$. The additive function $\omega_y(n)=\sum_{p\le y}\mathbf 1_{p\mid n}$ is a function of that vector. Under the model, the sum of independent Bernoulli random variables of means $1/p$ has mean and variance $\sum_{p\le y}1/p=\log\log y+O(1)=\tfrac12 L+O(1)$, and the classical central limit theorem, which Volume~4 proves for uniformly bounded independent summands with variance tending to infinity, yields a Gaussian after normalisation by $\sqrt{\log\log y}$.

The primes between $y$ and $x$ change $\omega$ by $\omega(n)-\omega_y(n)\le\log x/\log y=L^{1/2}$. Tur\'an--Kubilius, already proved in this chapter, bounds the second moment of $\omega-\omega_y$ on $n\le x$ by $O(L^{1/2})$. Chebyshev's inequality makes $(\omega-\omega_y)/\sqrt L\to0$ in probability. The means $\log\log x$ and $\log\log y$ differ by $\tfrac12 L$, which is the same order as the variance; replace $y$ by $x^{1/\log L}$, so that $\log\log y=L-\log\log L$ and $u=\log L\to\infty$, which is still inside the range of Theorem~\ref{thm:kubilius}. Then the truncated mean is $L+O(\log\log L)$, the Gaussian normalisation differs from $\sqrt L$ by $o(\sqrt L)$, and adding a quantity that is $o_P(\sqrt L)$ does not change the limiting distribution function. This is the corollary.
\end{proof}

Volume~6 proves the same limit by the method of moments, truncating at $z=x^{1/L^{1/4}}$ and comparing the moments of $\omega_z$ with those of the independent sum, then smoothing. The two arguments use the same truncation. They do not use the same bookkeeping.

What Volume~6 makes visible is the moment comparison: after the large prime factors are cut away, $\omega$ has the moments of a sum of independent coin tosses, and those moments are the Gaussian moments. The proof never builds a probability space on the primes; the independence is an asymptotic statement about averages of $\mathbf 1_{p\mid n}$.

What this chapter makes visible is the space. Theorem~\ref{thm:kubilius} is a total-variation comparison, so the coin tosses are a measure, not only a list of moments. The bell curve is the central limit theorem on that product, which Volume~4 already contains. The arithmetic content is the estimate that says the product is close to the integers.

The rate in the corollary is not the one R\'enyi and Tur\'an obtained from the Selberg--Delange formula. That rate is $O(1/\sqrt{\log\log x})$ in the Kolmogorov distance, and it uses the mean of $z^{\omega(n)}$ for complex $z$ of modulus $1$, which this chapter has proved and which Volume~6 does not need. Recording the rate would be a second pass through the same generating function, not a second foundation. Volume~11's description of a uniform integer modulo $\varphi$, which quotes the Selberg--Delange theorem from this chapter, is a consumer of the mean-value theorem, not of a second proof of Erd\H os--Kac.

\begin{proposition}[a Chernoff bound]\label{prop:chernoffomega}
Let $Q^*(\alpha)=\alpha\log\alpha-\alpha+1$ for $\alpha>0$, and fix $A\ge1$. Uniformly for $x\ge16$,
\[
P_x\bigl(\omega(n)\ge\alpha\log\log x\bigr)\ll_A(\log x)^{-Q^*(\alpha)}\qquad(1\le\alpha\le A).
\]
\end{proposition}
\begin{proof}
Selberg's formula in this chapter gives $\mathbb E_x z^{\omega(n)}\ll_A(\log x)^{z-1}$ for real $z\in[1,A]$. Markov's inequality at $z=\alpha$ produces the exponent $\alpha-1-\alpha\log\alpha=-Q^*(\alpha)$.
\end{proof}


\section{Exercises}

\begin{exercise}[$\star$]\label{ex:meanphi}
Show that $\sum_{n\le x}\frac{\varphi(n)}n=\frac6{\pi^2}x+O(\log x)$, using $\frac{\varphi(n)}n=\sum_{d\mid n}\frac{\mu(d)}d$. Check that $\frac6{\pi^2}$ is the mean of the limit law, $\prod_p\mathbb E[(1-\frac1p)^{\mathbf 1_{G_p\ge1}}]$.
\end{exercise}

\begin{exercise}[$\star$]\label{ex:omeganolimit}
Show that $\omega$ has no limiting distribution: for every $K$, $P_x(\omega(n)\le K)\to0$. Which of the three series diverges?
\end{exercise}

\begin{exercise}[$\star$]\label{ex:abundant6}
Show that every multiple of an abundant number is abundant, that $6k$ is abundant for $k\ge2$, and deduce that the density of abundant numbers is at least $\frac16+\frac1{20}-\frac1{60}=0.2$. (Use also $20$.)
\end{exercise}

\begin{exercise}[$\star\star$]\label{ex:phimoments}
Show that for every real $s>0$, $\frac1x\sum_{n\le x}(\frac{\varphi(n)}n)^s\to\prod_p\bigl(1-\frac1p+\frac1p(1-\frac1p)^s\bigr)$. Compute the limit for $s=1$ and $s=2$ as products, and evaluate them numerically.
\end{exercise}

\begin{exercise}[$\star\star$]\label{ex:purity}
Let $\varepsilon_k$ be independent with $P(\varepsilon_k=0)=P(\varepsilon_k=1)=\frac12$. Show that $S=\sum_{k\ge1}\varepsilon_k2^{-k}$ is uniform on $[0,1]$, and that $T=\sum_{k\ge1}2\varepsilon_k3^{-k}$ is continuous and carried by the Cantor set, hence purely singular. Check both against Theorems~\ref{thm:levyatoms} and~\ref{thm:jessenwintner}.
\end{exercise}

\begin{exercise}[$\star\star$]\label{ex:losharp}
Show that Lemma~\ref{lem:lo} is sharp up to a constant: for $a_1=\dots=a_m=1$, $\max_cP(\sum\varepsilon_ia_i=c)\sim\sqrt{\frac2{\pi m}}$. Show that if $a_1<\dots<a_m$ are distinct positive integers and $m\ge2$, then $P(\sum\varepsilon_ia_i=c)\le\frac12$, and compute this probability for $a=(1,2,3,4)$ and $c=0$.
\end{exercise}

\begin{exercise}[$\star\star$]\label{ex:kubiliusomega}
Use Corollary~\ref{cor:kubiliusf} to show that for $y=x^{1/\log\log x}$, the law of $\omega_y(n)=\#\{p\le y:p\mid n\}$ is within $o(1)$ in total variation of the law of a sum of independent Bernoulli variables with parameters $\frac1p$, $p\le y$. Deduce from the Chen--Stein bound of Volume~11 (Theorem~6.8) that $\omega_y(n)$ is within $O(\frac1{\log\log x})$ of a Poisson law of parameter $\sum_{p\le y}\frac1p$, up to the error of Theorem~\ref{thm:kubilius}.
\end{exercise}

\begin{exercise}[$\star\star\star$]\label{ex:phisupport}
Show that the support of the limit law of $\frac{\varphi(n)}n$ is the whole interval $[0,1]$: for every $0\le a<b\le1$, a positive proportion of integers have $a<\frac{\varphi(n)}n<b$. (Approximate a target $t\in\,]-\infty,0[$ by a finite sum of distinct terms $\log(1-\frac1p)$, using $\sum\frac1p=\infty$.)
\end{exercise}

\begin{exercise}[$\star$]\label{ex:hankelint}
Check Lemma~\ref{lem:hankel} directly for $z=1,2,3$ by the residue theorem, and for $z=0,-1$ by Cauchy's theorem. Use it to compute $\frac1{2\pi i}\int_{\mathcal H}w^{-1/2}e^w\dd w$.
\end{exercise}

\begin{exercise}[$\star$]\label{ex:twoomega}
Deduce from Corollary~\ref{cor:omegasd} that $\sum_{n\le x}2^{\omega(n)}=\frac6{\pi^2}x\log x+O(x)$. Then prove it without complex analysis, from $2^{\omega(n)}=\sum_{d\mid n}\mu(d)^2$.
\end{exercise}

\begin{exercise}[$\star\star$]\label{ex:minusone}
Show that $\sum_{n\le x}(-1)^{\omega(n)}=O(\frac x{(\log x)^3})$. Which value of the Gamma function is responsible? Numerically, one finds $\frac1x\sum_{n\le x}(-1)^{\omega(n)}=-0.000165$ at $x=10^7$.
\end{exercise}

\begin{exercise}[$\star\star$]\label{ex:bigomegasd}
Show that $\sum_{n\le x}2^{\Omega(n)}\ge c\,x(\log x)^2$ for $x\ge2$ and some $c>0$, so that Corollary~\ref{cor:bigomegasd} cannot extend to $z=2$. Where does its proof break down?
\end{exercise}

\begin{exercise}[$\star\star$]\label{ex:onlythree}
Show that the number of $n\le x$ all of whose prime factors are $\equiv3\pmod4$ is $\frac{c_3x}{\sqrt{\log x}}(1+O(\frac1{\log x}))$, with $c_3=\frac{\sqrt2}\pi\prod_{p\equiv3\ (4)}(1-\frac1{p^2})^{-1/2}=0.48652\ldots$
\end{exercise}

\begin{exercise}[$\star\star\star$]\label{ex:landaupik}
(Landau, 1900.) Let $\pi_k(x)=\#\{n\le x:\omega(n)=k\}$. Show, for each fixed $k\ge1$, that
\[
\pi_k(x)=\frac x{\log x}\,\frac{(\log\log x)^{k-1}}{(k-1)!}\Bigl(1+O_k\Bigl(\frac1{\log\log x}\Bigr)\Bigr).
\]
\end{exercise}

\section{Hints}

\begin{hint}{ex:meanphi}
Swap the sums: $\sum_{n\le x}\sum_{d\mid n}\frac{\mu(d)}d=\sum_{d\le x}\frac{\mu(d)}d\lfloor\frac xd\rfloor$.
\end{hint}
\begin{hint}{ex:omeganolimit}
Tur\'an--Kubilius with Chebyshev's inequality.
\end{hint}
\begin{hint}{ex:abundant6}
$\frac{\sigma(n)}n=\sum_{d\mid n}\frac1d$ only grows when $n$ gets new divisors.
\end{hint}
\begin{hint}{ex:phimoments}
Bounded continuous functions of a variable that converges in law.
\end{hint}
\begin{hint}{ex:purity}
Binary and ternary expansions; for $T$, count the intervals of the $k$-th stage of the Cantor construction.
\end{hint}
\begin{hint}{ex:losharp}
Stirling's formula; for distinct integers, look at what happens when only the largest $a_i$ changes sign.
\end{hint}
\begin{hint}{ex:kubiliusomega}
$\omega_y(n)=\Lambda(v_y(n))$ with $\Lambda(\mathbf k)=\#\{p:k_p\ge1\}$, and $\mathbf 1_{G_p\ge1}$ is Bernoulli$(\frac1p)$.
\end{hint}
\begin{hint}{ex:phisupport}
A finite set of primes $p_1<\dots<p_r$, all larger than some $P$, and the integers $n$ with $p_i\mid n$ and no other prime factor in $[2,y]$ except those above $P'$.
\end{hint}

\begin{hint}{ex:hankelint}
For $z=k\ge1$ the two edges cancel, since $w^{-k}$ is single-valued. For $z=\frac12$, $\Gamma(\frac12)=\sqrt\pi$.
\end{hint}
\begin{hint}{ex:twoomega}
$F(2)=\frac6{\pi^2}$ and $\Gamma(2)=1$. For the elementary proof, swap sums and use $\sum_{d\le y}\mu(d)^2=\frac6{\pi^2}y+O(\sqrt y)$.
\end{hint}
\begin{hint}{ex:minusone}
$\frac1{\Gamma(-1)}=0$.
\end{hint}
\begin{hint}{ex:bigomegasd}
Write $n=2^km$ with $m$ odd, and use $2^{\Omega(m)}\ge2^{\omega(m)}$ and Theorem~\ref{thm:sd} for the odd $m$.
\end{hint}
\begin{hint}{ex:onlythree}
Find $G$ with $F=\zeta^{1/2}G$; this time $L(s,\chi_4)$ appears with the exponent $-\frac12$.
\end{hint}
\begin{hint}{ex:landaupik}
Write
\[
\pi_k(x)=\frac1{2\pi i}\oint_{|z|=1/2}z^{-k-1}\sum_{n\le x}z^{\omega(n)}\dd z;
\]
the function $\Phi(z)=\frac{F(z)}{\Gamma(z)}$ is entire with $\Phi(z)=z+O(z^2)$.
\end{hint}

\section{Solutions}

\begin{solution}{ex:meanphi}
\idea Whenever an arithmetic function is a divisor sum $\sum_{d\mid n}g(d)$, its mean is $\sum_d\frac{g(d)}d$, by swapping the order of summation; the error is controlled by $\sum|g(d)|$.
\writeup Since $\varphi(n)=\sum_{d\mid n}\mu(d)\frac nd$, we have $\frac{\varphi(n)}n=\sum_{d\mid n}\frac{\mu(d)}d$, and
\[
\sum_{n\le x}\frac{\varphi(n)}n=\sum_{d\le x}\frac{\mu(d)}d\Bigl(\frac xd+O(1)\Bigr)=x\sum_{d\ge1}\frac{\mu(d)}{d^2}+O\Bigl(x\sum_{d>x}\frac1{d^2}\Bigr)+O\Bigl(\sum_{d\le x}\frac1d\Bigr)=\frac x{\zeta(2)}+O(\log x).
\]
In the model, $\frac{\varphi(n)}n$ becomes $\prod_p(1-\frac1p)^{\mathbf 1_{G_p\ge1}}$, a product of independent factors with mean $1-\frac1p+\frac1p(1-\frac1p)=1-\frac1{p^2}$. The product of the means is $\prod_p(1-p^{-2})=\frac1{\zeta(2)}=\frac6{\pi^2}$, and since $\frac{\varphi(n)}n$ is bounded the mean of the limit law is the limit of the means.
\end{solution}

\begin{solution}{ex:omeganolimit}
\idea If a function grows to infinity in probability, it cannot converge in law; Tur\'an--Kubilius shows that $\omega(n)$ is within $\sqrt{\log\log x}$ of $\log\log x$ for most $n$.
\writeup Here $f(p)=1$, so $f^*(p)=1$ and the second and third series are $\sum\frac1p=\infty$. With $A(x)=\sum_{p\le x}\frac1p=\log\log x+O(1)$ and $B(x)^2=A(x)$, Volume~6, Theorem~9.1, and Chebyshev give, for $x$ so large that $A(x)>2K$,
\[
P_x(\omega(n)\le K)\le P_x\Bigl(|\omega(n)-A|\ge\frac A2\Bigr)\le\frac{4C_0B^2}{A^2}(1+o(1))=\frac{4C_0}{A}(1+o(1))\to0.
\]
\end{solution}

\begin{solution}{ex:abundant6}
\idea Abundance only depends on the set of divisors, and multiplying $n$ adds divisors.
\writeup $\frac{\sigma(n)}n=\sum_{d\mid n}\frac1d$. If $n\mid N$, then $\frac{\sigma(N)}N=\sum_{e\mid N}\frac1e\ge\sum_{d\mid n}\frac1d=\frac{\sigma(n)}n$, since the divisors of $n$ are divisors of $N$. So multiples of abundant numbers are abundant. For $k\ge2$, $6k$ has the divisors $6k,3k,2k,k$ and $1$, so $\sigma(6k)\ge(6+3+2+1)k+1>12k$. The number $20$ has $\sigma(20)=42>40$. The multiples of $6$ or $20$ have density $\frac16+\frac1{20}-\frac1{\operatorname{lcm}(6,20)}=\frac16+\frac1{20}-\frac1{60}=\frac15$, and all but $6$ itself (which is perfect, $\sigma(6)=12$, hence abundant in our convention $\sigma(n)\ge2n$) are abundant. By Corollary~\ref{cor:phicont} the density exists, so it is at least $0.2$.
\end{solution}

\begin{solution}{ex:phimoments}
\idea Convergence in law means convergence of $\mathbb E\phi$ for bounded continuous $\phi$; $(\frac{\varphi(n)}n)^s$ is such a function of $\log\frac{\varphi(n)}n$, and the limit is computed in the independent model.
\writeup Let $\phi(t)=e^{st}$ for $t\le0$; it is bounded and continuous on $]-\infty,0]$, where $\log\frac{\varphi(n)}n$ lives. By Theorem~\ref{thm:erdoswintner}, $\mathbb E_x\phi(f(n))\to\mathbb E\phi(Z)$, and $e^{sZ}=\prod_p(1-\frac1p)^{s\mathbf 1_{G_p\ge1}}$ (strong additivity: only whether $G_p\ge1$ matters). The partial products converge almost surely and are bounded by $1$, so by dominated convergence and independence $\mathbb Ee^{sZ}=\prod_p(1-\frac1p+\frac1p(1-\frac1p)^s)$. For $s=1$ this is $\prod_p(1-p^{-2})=0.6079$. For $s=2$, $1-\frac1p+\frac1p(1-\frac2p+\frac1{p^2})=1-\frac2{p^2}+\frac1{p^3}$, and $\prod_p(1-\frac2{p^2}+\frac1{p^3})=0.4282$ (a numerical calculation evaluates the product with an Euler-product acceleration).
\end{solution}

\begin{solution}{ex:purity}
\idea Both are random digit expansions; the base and the allowed digits decide whether every interval gets its fair share.
\writeup For $S$, the binary digits of a uniform variable on $[0,1]$ are independent fair bits (Volume~4), and conversely the law of $S$ gives each dyadic interval $[\frac j{2^k},\frac{j+1}{2^k}]$ mass $2^{-k}$, which determines the uniform law. For $T$, the digits in base $3$ are $0$ or $2$, so $T$ lies in the Cantor set $C$, which at stage $k$ is covered by $2^k$ intervals of length $3^{-k}$: its Lebesgue measure is $0$. Each term $2\varepsilon_k3^{-k}$ has $q_k=\frac12$ and $\sum(1-q_k)=\infty$, so $T$ has no atom by Theorem~\ref{thm:levyatoms}; a law without atoms carried by a null set is purely singular. Both are pure, as Theorem~\ref{thm:jessenwintner} predicts; the two series differ only in the ratio $2$ versus $3$, which shows that the type cannot be read off from L\'evy's criterion alone.
\end{solution}

\begin{solution}{ex:losharp}
\idea For equal coefficients the sum is a centred binomial, whose largest probability is the central binomial coefficient. For distinct coefficients, flipping one sign moves the sum, so the sign vectors can be paired off.
\writeup With all $a_i=1$, $\sum\varepsilon_i=2|I|-m$, and the most likely value has probability $\binom m{\lfloor m/2\rfloor}2^{-m}$, which is $\sim\sqrt{\frac2{\pi m}}$ by Stirling's formula. For distinct positive $a_i$, pair each sign vector with the one in which only $\varepsilon_m$ is flipped. The two sums differ by $2a_m\ne0$, so at most one of each pair hits $c$, and $P\le\frac12$ (this uses only $a_m\ne0$). For $a=(1,2,3,4)$ and $c=0$, the set $I$ of plus signs must have $\sum_{I}a_i=5$, so $I=\{1,4\}$ or $I=\{2,3\}$, and $P=\frac2{16}=\frac18$. In fact for distinct coefficients the truth is of order $m^{-3/2}$, a theorem of S\'ark\"ozy and Szemer\'edi (1965).
\end{solution}

\begin{solution}{ex:kubiliusomega}
\idea $\omega_y(n)$ is a function of the exponent vector, so it inherits the total variation bound; in the model it is a sum of independent indicators, for which Volume~11 has a Poisson approximation.
\writeup $\omega_y(n)=\Lambda(v_y(n))$ with $\Lambda(\mathbf k)=\#\{p\le y:k_p\ge1\}$, and $\Lambda(G^{(y)})=\sum_{p\le y}\mathbf 1_{G_p\ge1}$, a sum of independent Bernoulli variables $B_p$ with $P(B_p=1)=P(G_p\ge1)=\frac1p$. By Theorem~\ref{thm:kubilius} with $u=\log\log x$ the two laws differ by $O(e^{-u/4}+x^{-1/4}\log x)=O((\log x)^{-1/4})$ in total variation. The Chen--Stein bound for a sum of independent Bernoulli variables gives $d_{TV}(\sum B_p,\operatorname{Poisson}(\lambda))\le\frac{1-e^{-\lambda}}\lambda\sum_p\frac1{p^2}$ with $\lambda=\sum_{p\le y}\frac1p=\log\log x+O(\log\log\log x)$; this is $O(\frac1{\log\log x})$. The triangle inequality gives the claim. The Erd\H{o}s--Kac section of this chapter shows that the full $\omega(n)$, which includes the large primes, is \emph{not} that close to a Poisson law: the large primes add a non-Poisson correction, the factor $\frac1{\Gamma(z)}$ of the Selberg--Delange theorem.
\end{solution}

\begin{solution}{ex:phisupport}
\idea To make $\frac{\varphi(n)}n$ close to a target, choose the set of primes dividing $n$; the divergence of $\sum\frac1p$ lets a greedy choice among large primes reach any target with any precision, and the Kubilius model shows that the integers with a prescribed set of small prime divisors have positive density.
\writeup Let $0<a<b\le1$ (the case $a=0$ follows by shrinking the interval), and fix $t$ with $\log a<t<\log b$, $t<0$. Let $\eta=\frac12\min(t-\log a,\log b-t)$ and choose $P$ with $|\log(1-\frac1p)|<\eta$ for $p>P$. Since $\sum_{p>P}|\log(1-\frac1p)|=\infty$ and each term is less than $\eta$, adding the terms one by one in increasing order of $p$ and stopping just before the total passes $|t|$ gives primes $P<p_1<\dots<p_r$ with $s=\sum_i\log(1-\frac1{p_i})\in\,]t-\eta,t]$. Next choose $Q>p_r$ so large that $\frac2\eta\sum_{p>Q}\frac1{p^2}<\frac12$. In the model, let $T=\sum_{p>Q}\log(1-\frac1p)\mathbf 1_{G_p\ge1}\le0$; since $|\log(1-\frac1p)|\le\frac2p$, Markov's inequality gives $P(|T|>\eta)\le\frac1\eta\sum_{p>Q}\frac2p\cdot\frac1p<\frac12$. The event $E$: ``$G_{p_i}\ge1$ for $i\le r$, $G_p=0$ for the other $p\le Q$, and $|T|\le\eta$'' has probability at least $\frac12\prod_i\frac1{p_i}\prod_{p\le Q}(1-\frac1p)>0$, by independence. On $E$, $\log Z=s+T\in[s-\eta,s]\subset\,]t-2\eta,t]\subseteq\,]\log a,\log b[$. So the limit law gives positive mass to $]a,b[$. By Theorem~\ref{thm:erdoswintner} and Corollary~\ref{cor:phicont} (the endpoints carry no mass), the proportion of $n\le x$ with $a<\frac{\varphi(n)}n<b$ tends to that positive mass.
\end{solution}

\section*{Sources and further reading}

\begin{itemize}
\item \emph{History.} I.~J.~Schoenberg, ``Über die asymptotische Verteilung reeller Zahlen mod~1'', \emph{Math.\ Z.}\ 28 (1928), 171--199, for $\frac{\varphi(n)}n$. H.~Davenport, ``Über numeri abundantes'', \emph{Sitzungsber.\ Preuss.\ Akad.\ Wiss.}\ (1933), 830--837. F.~Behrend, ``Über numeri abundantes II'', same \emph{Sitzungsber.}\ (1933), 280--293, for the bounds $0.241$ and $0.314$. M.~Deléglise, ``Bounds for the density of abundant integers'', \emph{Experiment.\ Math.}\ 7 (1998), 137--143; G.~McNew and S.~Setty, ``On the densities of covering numbers and abundant numbers'', arXiv:2507.23041 (2025), for the value $0.24761\ldots$
\item \emph{Erd\H{o}s--Wintner.} P.~Erd\H{o}s and A.~Wintner, ``Additive arithmetical functions and statistical independence'', \emph{Amer.\ J.\ Math.}\ 61 (1939), 713--721. P.~Erd\H{o}s, ``On the smoothness of the asymptotic distribution of additive arithmetical functions'', \emph{Amer.\ J.\ Math.}\ 61 (1939), 722--725.
\item \emph{Kubilius.} J.~Kubilius, \emph{Probabilistic Methods in the Theory of Numbers}, Translations of Mathematical Monographs~11, American Mathematical Society, 1964. Biographical facts from E.~Manstavičius, ``Jonas Kubilius (1921--2011)'', \emph{Acta Arith.}\ 157 (2013), 11--36, and the entry ``Jonas Kubilius'' of the \emph{Visuotinė lietuvių enciklopedija} (rector 1958--1990). P.~D.~T.~A.~Elliott, ``Jonas Kubilius and probabilistic number theory: some personal reflections'', \emph{Lith.\ Math.\ J.}\ 55 (2015), 2--24.
\item \emph{Probability.} P.~L\'evy, ``Sur les séries dont les termes sont des variables éventuelles indépendantes'', \emph{Studia Math.}\ 3 (1931), 119--155. B.~Jessen and A.~Wintner, ``Distribution functions and the Riemann zeta function'', \emph{Trans.\ Amer.\ Math.\ Soc.}\ 38 (1935), 48--88. P.~Erd\H{o}s, ``On a lemma of Littlewood and Offord'', \emph{Bull.\ Amer.\ Math.\ Soc.}\ 51 (1945), 898--902.
\item \emph{Textbooks.} G.~Tenenbaum, \emph{Introduction to Analytic and Probabilistic Number Theory}, 3rd ed., AMS, 2015, Chapters~III.2--III.4. P.~D.~T.~A.~Elliott, \emph{Probabilistic Number Theory~I}, Springer, 1979. E.~Kowalski, \emph{An Introduction to Probabilistic Number Theory}, Cambridge University Press, 2021, Chapter~2, whose presentation of the Kubilius model we follow.
\end{itemize}

\begin{solution}{ex:hankelint}
\idea Where $w^{-z}$ is single-valued, the loop can be closed into a circle and Hankel's formula becomes a residue computation.
\writeup For an integer $z$, $w^{-z}$ is single-valued, so the integrals along the two edges are equal and opposite and $\frac1{2\pi i}\int_{\mathcal H}w^{-z}e^w\dd w=\frac1{2\pi i}\oint_{|w|=1}w^{-z}e^w\dd w$. For $z=k\ge1$ this is the coefficient of $w^{k-1}$ in $e^w$, namely $\frac1{(k-1)!}=\frac1{\Gamma(k)}$: $1,1,\frac12$ for $k=1,2,3$. For $z=0$ or $-1$ the integrand $e^w$ or $we^w$ is entire and the integral vanishes, as does $\frac1\Gamma$ at $0$ and $-1$. For $z=\frac12$ the lemma gives $\frac1{\Gamma(1/2)}=\frac1{\sqrt\pi}$.
\end{solution}

\begin{solution}{ex:twoomega}
\idea Both proofs see the same constant: $\frac6{\pi^2}$ is the density of squarefree numbers, and $2^{\omega(n)}$ counts the squarefree divisors of $n$.
\writeup With $z=2$, $F(2)=\prod_p(1+\frac1p)(1-\frac1p)=\prod_p(1-\frac1{p^2})=\frac6{\pi^2}$ and $\Gamma(2)=1$, so Corollary~\ref{cor:omegasd} gives $x\log x(\frac6{\pi^2}+O(\frac1{\log x}))=\frac6{\pi^2}x\log x+O(x)$. Elementarily, $2^{\omega(n)}$ is the number of squarefree $d\mid n$, so $\sum_{n\le x}2^{\omega(n)}=\sum_{d\le x}\mu(d)^2\lfloor\frac xd\rfloor=x\sum_{d\le x}\frac{\mu(d)^2}d+O(x)$. By partial summation from $Q(y)=\sum_{d\le y}\mu(d)^2=\frac6{\pi^2}y+O(\sqrt y)$, $\sum_{d\le x}\frac{\mu(d)^2}d=\frac{Q(x)}x+\int_1^x\frac{Q(y)}{y^2}\dd y=\frac6{\pi^2}\log x+O(1)$.
\end{solution}

\begin{solution}{ex:minusone}
\idea At $z=-1$ the main term of the Selberg--Delange formula vanishes because $\frac1\Gamma$ has a zero at $-1$; only the error term remains.
\writeup Corollary~\ref{cor:omegasd} with $z=-1$, $R=1$: $\sum_{n\le x}(-1)^{\omega(n)}=x(\log x)^{-2}(\frac{F(-1)}{\Gamma(-1)}+O(\frac1{\log x}))$, and $\frac1{\Gamma(-1)}=0$, so the sum is $O(\frac x{(\log x)^3})$. In the language of the proof: $\zeta(s)^{-1}$ has no branch cut, the loop closes up, and the integrand $\zeta(s)^{-1}G(s)\frac{x^s}s$ is holomorphic at $s=1$. Delange's expansion of the Remark says more: every term contains $\frac1{\Gamma(-1-j)}=0$, so the sum is $O(\frac x{(\log x)^A})$ for every $A$. To do better one must use the zero-free region itself, as for the prime number theorem.
\end{solution}

\begin{solution}{ex:bigomegasd}
\idea The prime $2$ alone already does the damage: the powers of $2$ carry a weight $2^k$ that compensates exactly for their rarity.
\writeup For odd $m$, $a_m=2^{\omega(m)}\mathbf 1_{m\text{ odd}}$ has the Dirichlet series $\prod_{p>2}(1+\frac2{p^s-1})=\zeta(s)^2G_2(s)\frac{2^s-1}{2^s+1}$, with $G_2$ as in Lemma~\ref{lem:eulerprod}; the last factor is holomorphic and bounded on $\sigma>\frac34$ (there $|2^s+1|\ge2^{3/4}-1$) and equals $\frac13$ at $s=1$. Theorem~\ref{thm:sd} ($z=2$, $q=2$) gives $\sum_{m\le y,\,m\text{ odd}}2^{\omega(m)}\sim\frac{F(2)}3y\log y=\frac2{\pi^2}y\log y$, so this sum is at least $c\,y\log2y$ for all $y\ge1$ (for small $y$ the term $m=1$ suffices). Now $2^{\Omega(2^km)}=2^k2^{\Omega(m)}\ge2^k2^{\omega(m)}$, and with $K=\lfloor\log_2x\rfloor$,
\[
\sum_{n\le x}2^{\Omega(n)}\ge\sum_{k=0}^K2^k\,c\,\frac x{2^k}\log\frac{2x}{2^k}\ge c\,x\sum_{k=0}^K(K+1-k)\log2\ge\frac c{2\log2}\,x(\log x)^2.
\]
The proof of Corollary~\ref{cor:bigomegasd} breaks twice at $z=2$: the Euler factor at $p=2$ becomes $(1-2^{1-s})^{-1}$, with a pole at $s=1$, so $G$ is not holomorphic there; and no $q>1$ makes $\sum_{n\le y}2^{q\Omega(n)}$ of size $y(\log y)^K$, since the powers of $2$ alone give $y^q$.
\end{solution}

\begin{solution}{ex:onlythree}
\idea Write the Euler product, compare it with $\zeta(s)/L(s,\chi_4)$, and take a square root.
\writeup The series is $F(s)=\prod_{p\equiv3\ (4)}(1-p^{-s})^{-1}$. Prime by prime, $\frac{\zeta(s)}{L(s,\chi_4)}$ has the factor $(1-2^{-s})^{-1}$ at $2$, $1$ at $p\equiv1$, and $\frac{1+p^{-s}}{1-p^{-s}}$ at $p\equiv3$. So $F^2=(1-2^{-s})\frac\zeta L\prod_{p\equiv3}(1-p^{-2s})^{-1}$, and comparing logarithms as in Corollary~\ref{cor:landau}, $F=\zeta^{1/2}G$ with $G=(1-2^{-s})^{1/2}L(s,\chi_4)^{-1/2}\prod_{p\equiv3}(1-p^{-2s})^{-1/2}$. Lemma~\ref{lem:logzeta} bounds $|L^{-1/2}|$ by $e^{A_2/2}(\log\tau)^{1/2}$, so Theorem~\ref{thm:sd} applies with $z=\frac12$, and the constant is $\frac{G(1)}{\sqrt\pi}=\frac1{\sqrt\pi}\cdot\frac1{\sqrt2}\cdot\frac2{\sqrt\pi}\prod_{p\equiv3}(1-p^{-2})^{-1/2}=\frac{\sqrt2}\pi\prod_{p\equiv3}(1-p^{-2})^{-1/2}$. Numerically, one finds the ratio of the count to $\frac{c_3x}{\sqrt{\log x}}$ to be $1.026$ at $x=10^7$.
\end{solution}

\begin{solution}{ex:landaupik}
\idea The Selberg--Delange formula holds uniformly on a circle around $z=0$, so Cauchy's formula extracts the coefficient of $z^k$; the main term is a Taylor coefficient of $e^{z\log\log x}\Phi(z)$.
\writeup Let $\lambda=\log\log x$ and $L=\log x$. Since $\sum_{n\le x}z^{\omega(n)}=\sum_k\pi_k(x)z^k$ is a polynomial,
\[
\pi_k(x)=\frac1{2\pi i}\oint_{|z|=1/2}\frac{\sum_{n\le x}z^{\omega(n)}}{z^{k+1}}\dd z=\frac xL\,\frac1{2\pi i}\oint_{|z|=1/2}\frac{e^{\lambda z}\Phi(z)}{z^{k+1}}\dd z+E,
\]
with $\Phi=\frac F\Gamma$ and, by Corollary~\ref{cor:omegasd} with $R=1$, $|E|\le\frac1{2\pi}\pi\,2^{k+1}\max_{|z|=1/2}|xL^{z-1}|\cdot O(\frac1L)=O_k(xL^{-3/2})$. The main integral is the coefficient of $z^k$ in $e^{\lambda z}\Phi(z)$. The function $\Phi$ is entire, with $\Phi(z)=\sum_{j\ge1}\phi_jz^j$, because $\frac1{\Gamma(z)}=z+O(z^2)$ and $F(0)=\prod_p(1-\frac1p)(1-\frac1p)^{-1}=1$; so $\phi_1=1$. The coefficient is $\sum_{j=1}^k\phi_j\frac{\lambda^{k-j}}{(k-j)!}=\frac{\lambda^{k-1}}{(k-1)!}(1+O_k(\frac1\lambda))$. Since $xL^{-3/2}=\frac xL\cdot L^{-1/2}$ is much smaller than $\frac xL\frac{\lambda^{k-2}}{(k-1)!}$, the claim follows. For $k=1$ this is the prime number theorem with $\pi_1(x)=\pi(x)+\#\{p^a\le x,a\ge2\}$, in the weak form $\frac x{\log x}(1+O(\frac1{\log\log x}))$.
\end{solution}

\section*{Sources and further reading}

\begin{itemize}
\item \emph{Landau and Ramanujan.} E.~Landau, ``Über die Einteilung der positiven ganzen Zahlen in vier Klassen nach der Mindestzahl der zu ihrer additiven Zusammensetzung erforderlichen Quadrate'', \emph{Arch.\ Math.\ Phys.}\ (3) 13 (1908), 305--312. B.~C.~Berndt and R.~A.~Rankin, \emph{Ramanujan: Letters and Commentary}, AMS/LMS, 1995, for the letter of 16 January 1913. G.~H.~Hardy, \emph{Ramanujan: Twelve Lectures on Subjects Suggested by His Life and Work}, Cambridge, 1940, pp.~9 and~63. D.~Shanks, ``The second-order term in the asymptotic expansion of $B(x)$'', \emph{Math.\ Comp.}\ 18 (1964), 75--86. P.~Moree, ``Counting numbers in multiplicative sets: Landau versus Ramanujan'', \emph{Mathematics Newsletter} 21 (2011), 73--81 (arXiv:1110.0708), for the story of Ramanujan's reply, of Stanley's computation and of Shanks's verdict; P.~Moree and J.~Cazaran, ``On a claim of Ramanujan in his first letter to Hardy'', \emph{Expo.\ Math.}\ 17 (1999), 289--312.
\item \emph{Sathe and Selberg.} L.~G.~Sathe, ``On a problem of Hardy on the distribution of integers having a given number of prime factors'', I--II, \emph{J.\ Indian Math.\ Soc.}\ 17 (1953), 63--141; III--IV, 18 (1954), 27--81. A.~Selberg, ``Note on a paper by L.~G.~Sathe'', \emph{J.\ Indian Math.\ Soc.}\ 18 (1954), 83--87; its opening paragraph calls Sathe's results very beautiful and his method complicated and involved.
\item \emph{Delange.} H.~Delange, ``Sur des formules dues à Atle Selberg'', \emph{Bull.\ Sci.\ Math.}\ (2) 83 (1959), 101--111; ``Sur des formules de Atle Selberg'', \emph{Acta Arith.}\ 19 (1971), 105--146.
\item \emph{Landau 1900.} E.~Landau, ``Sur quelques problèmes relatifs à la distribution des nombres premiers'', \emph{Bull.\ Soc.\ Math.\ France} 28 (1900), 25--38, for $\pi_k(x)$.
\item \emph{Hankel.} H.~Hankel, ``Die Euler'schen Integrale bei unbeschränkter Variabilität des Arguments'', \emph{Z.\ Math.\ Phys.}\ 9 (1864), 1--21.
\item \emph{Textbooks.} G.~Tenenbaum, \emph{Introduction to Analytic and Probabilistic Number Theory}, 3rd ed., AMS, 2015, Chapters~II.5--II.6, whose formulation of the method we follow. H.~L.~Montgomery and R.~C.~Vaughan, \emph{Multiplicative Number Theory~I}, Cambridge, 2007, Chapter~7. R.~de~la~Bretèche and G.~Tenenbaum, ``Remarks on the Selberg--Delange method'', \emph{Acta Arith.}\ 200 (2021), 349--369 (arXiv:2010.12929).
\end{itemize}
